===== BEGIN FILE: TR-ARPES-X.tex =====
\documentclass[twocolumn,aps,a4paper,superscriptaddress]{revtex4-2}
\usepackage{kpfonts}
\usepackage{amsmath}
\usepackage{amssymb}
\usepackage{times}
\usepackage{graphicx}
\usepackage[colorlinks,linkcolor=blue,citecolor=blue,urlcolor=blue,breaklinks]{hyperref}
\usepackage[normalem]{ulem}
\input{ourmacros}

\begin{document}

\title{Unified First-Principles Theory of Coherent and Incoherent Excitons in Time-Resolved ARPES}
\title{Unified First-Principles Formula for Time-Resolved ARPES
Spectra of Coherent and Incoherent Excitons}
\author{Gianluca Stefanucci}
\affiliation{Dipartimento di Fisica, Universit{\`a} di Roma Tor Vergata, Via della Ricerca Scientifica 1,
00133 Rome, Italy}
\affiliation{INFN, Sezione di Roma Tor Vergata, Via della Ricerca Scientifica 1, 00133 Rome, Italy}

\author{Enrico Perfetto}
\affiliation{Dipartimento di Fisica, Universit{\`a} di Roma Tor Vergata, Via della Ricerca Scientifica 1,
00133 Rome, Italy}
\affiliation{INFN, Sezione di Roma Tor Vergata, Via della Ricerca Scientifica 1, 00133 Rome, Italy}




\begin{abstract}
Despite major experimental progresses in time-resolved
and angle-resolved photoemission spectroscopy, a  quantitative,
microscopic framework for interpreting
exciton-induced modifications of electronic band structures --
applicable even beyond the low-density limit -- is still lacking.
Here we close this gap by introducing a unified
approach that links the dynamics of
coherent and incoherent excitons to distinct and experimentally observable
excitonic sidebands.
Our central result is a general, first-principles
formula for time-resolved photoemission spectra, applicable across a
broad range of temperatures, excitation densities, and pump–probe
delays. This advance provides a predictive tool for
quantitatively tracking excitonic dynamics in complex
materials.
\end{abstract}

\maketitle




{\em Introduction.--}
The concept of excitons as quasiparticles in semiconductors and
insulators has long been central to the interpretation of optical
spectra, where it underpins the understanding of subgap resonances
and above-gap structures in absorbance and reflectance
measurements~\cite{onida_electronic_2002}. Optical probes, however,
are fundamentally restricted to bright, zero–center-of-mass-momentum
excitons and primarily access their energies. By contrast,
angle-resolved photoemission spectroscopy (ARPES) reveals
a far
richer and more detailed
picture~\cite{zhang_angle-resolved_2022,boschini_time-resolved_2024,xu_time-domain_2025}.
In particular, time-resolved ARPES (TR-ARPES)
has recently
emerged as a powerful platform for exciton physics, enabling  direct access
to bright, momentum-dark, and spin-dark
excitons~\cite{madeo_directly_2020,dendzik_observation_2020,tanimura_time_2020,wallauer_momentum-resolved_2021,fukutani_detecting_2021,kunin_momentum-resolved_2023,bange_probing_2024,gosetti_unveiling_2024}.
Beyond exciton energies, TR-ARPES uniquely resolves excitonic
wavefunctions in both band and momentum
space~\cite{man_experimental_2021,dong_direct_2021,meneghini_hybrid_2023,mori_spin-polarized_2023}.

Although TR-ARPES~\cite{na_advancing_2023} has led to
informative experimental spectra,  a rigourous
connection with the underlying exciton dynamics  is still
lacking.
After pumping below the bandgap,
the nonequilibrium exciton
fluid~\cite{haug2009quantum,lindberg_effective_1988,ostreich_non-stationary_1993,hannewald_excitonic_2000,koch_semiconductor_2006}
drives the motion of the nuclear
lattice.
The scattering between electrons and phonons
destroys exciton
coherence~\cite{madrid_phonon-induced_2009,nie_ultrafast_2014,dey_optical_2016,sangalli_an-ab-initio_2018,stefanucci_from-carriers_2021,perfetto_real-time_2023}
and diffuses
excitons~\cite{bertoni_generation_2016,uddin_neutral_2020,perfetto_ultrafast_2021,stefanucci_excitonic_2025}.
Bright coherent excitons are converted
-- typically in less than a few hundreds of femtoseconds~\cite{moody_intrinsic_2015,selig_exciton_2016,jakubczyk_impact_2018} -- into
bright and dark incoherent excitons.
TR-ARPES spectra of {\em coherent} excitons have been successfully
described~\cite{perfetto_pump-driven_2019,perfetto_time-resolved_2020,perfetto_floquet_2020,chan_giant_2023,pareek_driving_2026}.
However,  a first principles many-body framework for
the {\em incoherent} regime -- applicable even beyond the
domain of low excitation densities -- remains underdeveloped~\cite{perfetto_first-principles_2016,steinhoff_exciton_2017,rustagi_photoemission_2018,christiansen_theory_2019}.
In this Letter, we identify distinct signatures of coherent and
incoherent excitons in the TR-ARPES signal and provide a unified
theory of excitonic sidebands emerging from different exciton states.
The main finding is a
ready-to-use first-principles formula for interpreting TR-ARPES spectra
across a wide range
of temperatures, excitation densities and pump-probe delays.

{\em Photocurrent.--}
The rate of photoelectrons with energy $\e$ and parallel
momentum $\blk$, ejected by a high-energy probe
pulse $a(t)=ae^{-i\w_{0}t}+{\rm h.c.}$
impinging on the crystal after a delay $\t$ from the pump, is
proportional to~\cite{freericks_theoretical_2009,perfetto_first-principles_2016}
\begin{align}
I_{\blk}(\t,\e)&\propto-i|a|^{2}\sum_{\m\m'}
D_{\m\blk}(\e)D^{\ast}_{\m'\blk}(\e)
G^{<}_{\m\m'\blk}(\t,\w_{0}-\e),
\label{tdph2}
\end{align}
where  $D_{\m\blk}(\e)$ is the
photoemission matrix element for an electron
with band-momentum index $(\m\blk)$.
Henceforth, we use $\m=c$ for conduction bands and $\m=v$ for valence
bands. The nonequilibrium electronic
properties of the crystal are encoded in the Fourier transform,
$G^{<}(\t,\w)$,
of the lesser Green's function $G^{<}(t,t')$ with respect to the
time-difference $t-t'$ for a given center-of-mass-time (delay) $\t=(t+t')/2$.
In the quasiparticle approximation
$G^{<}_{\m\m'\blk}(\t,\w)=\d_{\m\m'}2\p i
f_{\m\blk}(\t)\d(\w-\e_{\m\blk}(\t))$, and the spectrum is peaked at
the quasiparticle energies $\e_{\m\blk}(\t)$ -- with an intensity
proportional to the
nonequilibrium populations $f_{\m\blk}(\t)$.
Since our primary interest is in the excitonic sidebands appearing
above the valence band maximum,
we  focus exclusively on the contribution of conduction electrons  to
the photoemission signal. Consequently, we restrict
the sum over $\m,\m'$ in Eq.~(\ref{tdph2}) to the conduction bands.

{\em TR-ARPES formula.--}
The interacting
nonequilibrium $G^{<}$ for conduction electrons can be calculated
from~\cite{svl-book_2025}
\begin{align}
G^{<}_{\blk}(\t,\w)=G^{R}_{\blk}(\t,\w)\S^{<}_{\blk}(\t,\w)G^{A}_{\blk}(\t,\w),
\label{G<}
\end{align}
where all quantities are matrices with conduction band indices.
We remark that this way of calculating $G^{<}$
differs from
earlier
attempts~\cite{perfetto_first-principles_2016,steinhoff_exciton_2017}, which
rely on the fluctuation-dissipation {\em ansatz}
$G^{<}_{\blk}(\t,\w)=i f(\w-\m_{c})A_{\blk}(\t,\w)$, where $f$ is the Fermi
function with a conduction-band
chemical potential $\m_{c}$ fixed by the excitation density,
and $A_{\blk}(\t,\w)$ is the nonequilibrium spectral function. This
ansatz is problematic for frequencies $\w$ below the gap since, at finite
temperature and/or out of
equilibrium, the addition and removal energies
overlap.
The retarded and advanced Green's functions
\begin{align}
G^{R}_{\blk}(\t,\w)=[G^{A}_{\blk}(\t,\w)]^{\dag}=[\w-\callE_{\blk}(\t)-\S^{R}_{\blk}(\t,\w)],
\end{align}
with $[\callE_{\blk}(\t)]_{\m\m'}\equiv \d_{\m\m'}\e_{\m\blk}(\t)$,
and the lesser
Green's function in Eq.~(\ref{G<}) can be calculated once an approximation for the
electronic self-energy $\S$ is available.


Let us consider  pump fields with frequencies
below the bandgap, ensuring that only discrete states are
explored during the time evolution.
We demonstrate below that Hartree plus statically
screened exchange (HSEX) effectively capture
the coherent regime, while T-matrix plus exchange
accurately describes the incoherent regime.
Remarkably, the full self-energy can be compactly
expressed as
\begin{subequations}
\begin{align}
\S^{R}_{cc'\blk}(\t,\w)&=\sum_{\l\blQ v}\s^{\l\blQ
v}_{cc'\blk}(\t,\w)\;,
\label{sigmaR}
\\
\S^{<}_{cc'\blk}(\t,\w)&=\sum_{\l\blQ v}
f(\w-E_{\l\blQ}(\t)-\e_{v\blk-\blQ}(\t))
\nn\\
&\times[\s^{\l\blQ v}_{cc'\blk}(\t,\w)-\s^{\l\blQ
v\ast}_{cc'\blk}(\t,\w)]\;,
\label{sigma<}
\end{align}
\label{sigma}
\end{subequations}
where the sum runs over all valence bands and over all exciton states $\l$ of
momentum $\blQ$. In the Fermi function of Eq.~(\ref{sigma<}),
the quasiparticle energies
$\e_{\m\blk}(\t)$ and exciton energies $E_{\l\blQ}(\t)$ are those of
the {\em nonequilibrium} system.
The key quantity in Eq.~(\ref{sigma}) is the function
\begin{align}
\s^{\l\blQ v}_{cc'\blk}(\t,\w)=K^{\l\blQ}_{cv\blk}\;
\frac{N_{\l\blQ}(\t)}{\w-E_{\l\blQ}(\t)-\e_{v\blk-\blQ}(\t)+i\eta}\;
K^{\l\blQ\ast}_{c'v\blk}-i\eta\;,
\label{sigmasmall}
\end{align}
where $\eta$ is a positive infinitesimal, and the amplitude
\begin{align}
K^{\l\blQ}_{cv\blk}=(\e^{\rm eq}_{c\blk}-\e^{\rm eq}_{v\blk-\blQ}-E^{\rm eq}_{\l\blQ})A^{\l
\blQ}_{cv\blk-\blQ}
\label{contrdef}
\end{align}
is expressed in terms of the
equilibrium band structure $\e^{\rm eq}_{\m\blk}$,
exciton energies $E^{\rm eq}_{\l\blQ}$ and wavefunctions $A^{\l
\blQ}_{cv\blk}$. Notably, the full self-energy
depends only on the {\em total} number of
excitons at the probing time $\t$~\cite{stefanucci_excitonic_2025}
\begin{align}
N_{\l\blQ}(\t)=\d_{\blQ,\bz}|\r_{\l}(\t)|^{2}+
N^{\rm inc}_{\l\blQ}(\t),
\label{N=c+ixb}
\end{align}
with $|\r_{\l}|^{2}$ the number of coherent excitons,
and $N^{\rm inc}_{\l\blQ}$ the number of incoherent
excitons.
In Appendix~\ref{findensapp} we present a
practical method to extract the nonequilibrium energies
$\e_{\m\blk}(\t)$ and $E_{\l\blQ}(\t)$
directly from $N_{\l\blQ}(\t)$.
Equations~(\ref{G<}-\ref{sigmasmall}) are the main
result of this Letter. They
 provide a first-principles formula
to calculate TR-ARPES spectra for any pump-probe delay, and for a
wide range of temperatures and excitation densities.
The only input quantities are
the exciton
populations $N_{\l\blQ}(\t)$, which
can be determined by, e.g., solving the excitonic Bloch
equations~\cite{thranhardt_quantum_2000,selig_exciton_2016,brem_exciton_2018,stefanucci_excitonic_2025}.


{\em Coherent regime.--} Here we briefly revisit the
theory of the coherent regime, and offer an alternative perspective.
Without any
loss of generality we set
the time origin, $t=0$, after the action of the pump pulse. For weak
pumping the excitation density is low, and
the state of the system at positive times is given by the quantum
superposition
$|\Q(t)\ket=|\Q_{g}\ket+\sum_{\l}\r_{\l}(t)|\l\bz\ket$, where
$|\Q_{g}\ket$ is the ground state and
$|\l\blQ\ket=\sum_{cv\blk}A^{\l\blQ}_{cv\blk}\hat{d}^{\dag}_{c\blk+\blQ}\hat{d}_{v\blk}|\Q_{g}\ket$
are the exciton states of momentum $\blQ$ -- here
$\hat{d}_{\m\blk}$ is the annihilation operator for electrons in band
$\m$ with momentum $\blk$. As the pump pulse
can only excite bright excitons,  $|\Q(t)\ket$ contains
exclusively excitons of vanishing momentum, $\blQ=\bz$. The coefficients
$\r_{\l}(t)=\sum_{cv\blk}A^{\l\bz\ast}_{cv\blk}\,
\bra\Q(t)|\hat{d}^{\dag}_{v\blk}\hat{d}_{c\blk}|\Q(t)\ket$ are the
exciton polarizations, and their square modulus yield the number of
coherent excitons, see Eq.~(\ref{N=c+ixb}). The
system is described by the state $|\Q(t)\ket$
up to times smaller than the exciton lifetimes.
Such time window defines the coherent regime. From
the knowledge of $|\Q(t)\ket$ we can directly calculate $G^{<}$, see
Appendix~\ref{lehmannapp}, and hence the
TR-ARPES spectrum. We refer to this strategy as the {\em Lehmann
approach}, since it is based on expanding the many-body state in
eigenstates of the Hamiltonian.
The same $G^{<}$ can be
obtained from Eq.~(\ref{G<}) using
the HSEX self-energy, see Eq.~(\ref{G<tmatfinld}).
The explicit form
of the HSEX self-energy can be readily derived
by projecting the equations of motion for the Green's
function~\cite{perfetto_pump-driven_2019,perfetto_time-resolved_2020} onto
the conduction-band subspace. The final result is
Eqs.~(\ref{sigma}) with $\s$ evaluated in
$N_{\l\blQ}=\d_{\blQ\bz}|\r_{\l}|^{2}$, see
Appendices~\ref{hsextwobandapp} and~\ref{hsexmultbandapp} for details.
The advantage of the many-body formulation over the
Lehmann approach is
that it successfully extends beyond the domain of low excitation
densities~\cite{pareek_driving_2026}.

{\em Incoherent regime.--}
In the incoherent regime the
polarization $\r_{\l}$ vanishes, and consequently the HSEX
self-energy in Eqs.~(\ref{sigma}) vanishes too.
In this regime, the
electronic subsystem is described by an ensemble of bright and dark
incoherent excitons.
The transition from a pure state to an ensemble is a general
phenomenon in systems interacting with a bath, such as
electrons interacting with phonons in our case.
It has been demonstrated in an exactly  solvable two-band
model~\cite{stefanucci_from-carriers_2021,zhang_long-lived_2025},
and it is also a consequence of the excitonic Bloch
equations~\cite{thranhardt_quantum_2000,selig_exciton_2016,brem_exciton_2018,stefanucci_excitonic_2025}.

A heuristic theory for the photocurrent in the incoherent regime
can be effectively developed for low excitation densities.
According to the excitonic Bloch equations the incoherent ensemble is described
by the many-body density matrix $\hat{\r}=|\Q_{g}\ket\bra\Q_{g}|+\sum_{\l\blQ}N^{\rm
inc}_{\l\blQ}|\l\blQ\ket\bra\l\blQ|$. From the knowledge of
$\hat{\r}$ we can calculate $G^{<}$, see Appendix~\ref{lehmannapp},
which in turn allows us to determine
the TR-ARPES spectrum. This is the Lehmann approach of the incoherent
regime~\cite{perfetto_first-principles_2016,sangalli_an-ab-initio_2018,rustagi_photoemission_2018,stefanucci_from-carriers_2021,chen_first-principles_2022}.
Establishing the many-body diagrammatic approximation
underlying this approach
would lift the low-density constraint, and pave the way toward
a fully first-principles treatment of TR-ARPES spectra in the
incoherent  regime.

At present, the state-of-the-art many-body
approximation is largely based on
the T-matrix diagrams~\cite{kremp_equation_1984,semkat_ionization_2009,perfetto_first-principles_2016,steinhoff_exciton_2017}.
Previous studies reported
excitonic sidebands appearing as replicas
of the conduction bands~\cite{perfetto_first-principles_2016,steinhoff_exciton_2017},
in apparent contradiction to the Lehmann
approach, which instead predicts replicas of the valence bands.
Another issue pertains to the intensity of the photoemission signal,
proportional to $|A^{\l\blQ}_{cv\blk}|^{2}$. The
exciton wavefunctions satisfy the Bethe
Salpeter equation (BSE), whose
kernel is the sum of the bare
Coulomb interaction $v$ (direct term) and the screened interaction $W$
(exchange term)~\cite{onida_electronic_2002,strinati_application_1988,svl-book_2025}, i.e.,
\begin{align}
(\e^{\rm eq}_{c\blk+\blQ}-\e^{\rm eq}_{v\blk})A^{\l\blQ}_{cv\blk}-
\sum_{c'v'\blk'}
K^{\blQ}_{cv\blk,c'v'\blk'}A^{\l\blQ}_{c'v'\blk'}=
E^{\rm eq}_{\l\blQ}A^{\l\blQ}_{cv\blk}\;,
\label{eigeneqX}
\end{align}
with
\begin{align}
K^{\blQ}_{cv\blk,c'v'\blk'}=
W_{c\blk+\blQ\,v'\blk'\,v\blk\,c'\blk'+\blQ}-
v_{c\blk+\blQ\,v'\blk'\,c'\blk'+\blQ\,v\blk}.
\label{hsex}
\end{align}
Thus, a pure T-matrix self-energy (containing only
$W$) cannot yield the
Lehmann intensity of the signal.

To reconcile many-body theory with the  Lehmann approach we consider the four self-energy
diagrams of Fig.~\ref{Tmatself}  -- the gray
circle is the exchange-correlation (xc) function $L$, the wiggly line
is the bare interaction $v$, and the double wiggly line is the
statically screened interaction $W$.
The first diagram in Fig.~\ref{Tmatself} is the T-matrix self-energy,
and  we refer to the sum of all four diagrams as the T-matrix plus
exchange (TX) self-energy.
Notice that the zeroth order xc function is simply $GG$. Consequently,
the lowest-order T-matrix diagram
has the same topology as the (second-Born) bubble diagram,
except that {\em both} interaction lines are screened.
This raises the issue of a possible double counting of screening
diagrams, which we address below.


The diagrammatic equation of Fig.~\ref{Tmatself} is
an approximation for the {\em lesser} and {\em greater} self-energy,
not for the time (or contour) ordered $\S$. Accordingly,
the bubble diagram with two $W$ lines does not emerge from T-matrix,
rather from GW. Indeed,
$\S_{\rm GW}^{<}=iG^{<}W^{<}=iG^{<}W^{R}P^{<}W^{A}$ where
$P^{<}=iG^{<}G^{>}$~\cite{svl-book_2025}.
Evaluating the screened interaction in
the static approximation, $W^{R}=W^{A}=W$, we recover
the lowest order (first)
diagram in Fig.~\ref{Tmatself}. Thus, there is no double counting
issue. We also observe that
the first diagram does not include
the {\em full} GW self-energy, since the bottom Green's function
is summed over
conduction bands only.
The missing contributions are recovered by adding the
fourth diagram in Fig.~\ref{Tmatself}. In fact, this diagram has the
structure $Gv\chi v=G(W-v)$, $\chi$ being the density response
function, and
$G$ is summed over the valence bands. In summary,  the
lesser and greater self-energies of Fig.~\ref{Tmatself} incorporate the full GW self-energy,
but handle $W$ differently depending on the Green's function band indices.
If $G$ has conduction indices then
$W$ is statically screened, whereas for valence indices $W$ is
dynamically screened.

\begin{figure}[t]
    \centering
\includegraphics[width=0.48\textwidth]{sigma.pdf}
\caption{TX self-energy for the Green's function in the incoherent
regime. We represent Green's functions $G$ with solid
lines, bare interactions $v$ with wiggly lines, screened interactions
$W$ with wiggly lines and the xc function $L$ with
a gray circle.}
\label{Tmatself}
\end{figure}

Taking into account the definition of the HSEX kernel in
Eq.~(\ref{hsex}), the TX self-energy of Fig.~\ref{Tmatself}
can be written as
\begin{align}
\S^{\lessgtr}_{cc'\blk}(t,t')&=-i^{2}
\sum_{\substack{c_{1}c_{2}v_{1}v_{2}v\\\blk_{1}\blk_{2}\blQ}}
K^{\blQ}_{cv\blk-\blQ,c_{1}v_{1}\blk_{1}}
L^{\blQ\lessgtr}_{c_{1}v_{1}\blk_{1},c_{2}v_{2}\blk_{2}}(t,t')
\nn\\ &\times
K^{\blQ\ast}_{c'v\blk-\blQ,c_{2}v_{2}\blk_{2}}G^{\lessgtr}_{v\blk-\blQ}(t,t'),
\label{sigma<>}
\end{align}
from which we can extract the retarded self-energy according to
\begin{align}
\S^{R}_{cc'\blk}(\t,\w)&=i\int\frac{d\w'}{2\p}
\frac{\S^{>}_{cc'\blk}(\t,\w)-\S^{<}_{cc'\blk}(\t,\w)}{\w-\w'+i\eta}-i\eta.
\label{sigmaRTX}
\end{align}
Expanding the xc function $L$ in the excitonic
basis~\cite{perfetto_first-principles_2016,stefanucci_excitonic_2025},
and approximating the
Green's function at the quasiparticle level,
we obtain Eq.~(\ref{sigmaR}) with $\s$ evaluated in $N_{\l\blQ}=N^{\rm
inc}_{\l\blQ}$ -- we refer the reader to
Appendix~\ref{txapp} for the derivation.
This result agrees with that of a pure T-matrix
calculation~\cite{steinhoff_exciton_2017} only if the exciton
energies and wavefunctions
remain unchanged upon setting
the bare $v$ to zero
in Eq.~(\ref{hsex}), which is generally not the case.

Equation~(\ref{sigma<>}) is only used to
extract $\S^{R}$.
The lesser self-energy in Eq.~(\ref{G<}) is derived from
 $\s$ according to Eq.~(\ref{sigma<}), ensuring consistency with the
self-energy in the coherent regime, see Appendix~\ref{hsexmultbandapp}.
As we see below [Eq.~(\ref{G<i2})], such consistency
is essential not only for recovering the results of the Lehmann
approach but also for extending the theory beyond the low-density
limit.

{\em Low density.--}
In the low-density limit  $N_{\l\blQ}$ is infinitesimal, and
$E_{\l\blQ}(\t)$ and $\e_{\m\blk}(\t)$
are well approximated by their equilibrium values.
To lowest order in the
excitation density, Eq.~(\ref{G<}) yields
\begin{align}
G^{<}_{cc'\blk}&(\t,\w)=\lim_{\eta\to 0}\frac{1}{\w-\e^{\rm eq}_{c\blk}+i\eta}
\S^{<}_{cc'\blk}(\t,\w)\frac{1}{\w-\e^{\rm eq}_{c'\blk}-i\eta}
\nn\\
=&2\p i\sum_{\l\blQ v}N_{\l\blQ}(\t)\;A^{\l\blQ}_{cv\blk-\blQ}
A^{\l\blQ\ast}_{c'v\blk-\blQ}\d(\w-E^{\rm eq}_{\l\blQ}-\e^{\rm eq}_{v\blk-\blQ}),
\label{G<tmatfinld}
\end{align}
where we use Eq.~(\ref{contrdef}).
This result agrees with the Lehmann approach both in the coherent
($N_{\l\blQ}=\d_{\blQ\bz}|\r_{\l}|^{2}$) and
incoherent ($N_{\l\blQ}=N_{\l\blQ}^{\rm inc}$) regimes, see
Appendix~\ref{lehmannapp}.
It also agrees with the
removal part of the spectral function  in the incoherent regime presented in
Ref.~\cite{steinhoff_exciton_2017}.

\begin{figure*}[tbp]
    \centering
\includegraphics[width=0.98\textwidth]{spectrafig.pdf}
\caption{(a) First Brillouin zone and location of a few high symmetry
points. (b) Exciton band structure of the equilibrium  WSe$_{2}$ ML
along the red path illustrated in panel~(a) -- curves are colored
according to the thermal occupation of excitons at temperature
$T=70$~K and excitation density $n_{c}=10^{13}$~cm$^{-2}$. (c,f)
Illustration of conduction band and replica of valence band at low
[panel~(c)] and moderate [panel~(f)]
excitation densities. (d,e,g,h) TR-ARPES spectrum (in arbitrary
units) for low [panels~(d,e)] and high [panels~(g,h)] excitation
densities at different temperatures:
$n_{c}=10^{11}$~cm$^{-2}$ and $T=0$~K [panel~(d)];
$n_{c}=10^{11}$~cm$^{-2}$ and  $T=70$~K
[panel~(e)]; $n_{c}=4\times 10^{12}$~cm$^{-2}$ and
$T=0$~K [panel~(g)];  $n_{c}=10^{13}$~cm$^{-2}$ and  $T=70$~K [panel~(h)].
}
\label{spectrafig}
\end{figure*}



Inserting Eq.~(\ref{G<tmatfinld}) in the photocurrent formula
Eq.~(\ref{tdph2}) we find
\begin{align}
I_{\blk}(\t,\e)&\propto |a|^{2}\sum_{\l\blQ}N_{\l\blQ}(\t)\sum_{v}
\Big|\sum_{c}D_{c\blk}(\e)A^{\l\blQ}_{cv\blk-\blQ}\Big|^{2}
\nn\\
&\times 2\p\d(\e^{\rm eq}_{v\blk-\blQ}+E^{\rm eq}_{\l\blQ}-\e+\w_{0}).
\label{tdphincoh1}
\end{align}
This formula extends the theory of
Refs.~\cite{perfetto_first-principles_2016,rustagi_photoemission_2018}
to multiple bands, and corrects the ansatz in
Ref.~\cite{chen_first-principles_2022} by placing the sum over
conduction bands inside the square modulus.
According to Eq.~(\ref{tdphincoh1}), a $\l\blQ$ exciton generates a
replica of the valence bands shifted in energy by $E^{\rm eq}_{\l\blQ}$
and in momentum by $\blQ$, see Fig.~\ref{spectrafig}(c).
It is worth observing that if the lowest-energy exciton, say $\l_{0}$,
is bright, and if we pump resonantly with it, then
$\r_{\l}(\t)\propto \d_{\l\l_{0}}$ (resonance condition).
As phonons do not scatter the exciton into higher energy states, the
coherent $\l_{0}$-exciton can only transition into an incoherent
$\l_{0}$-exciton of vanishing momentum. Considering the conservation of the total number of
excitons~\cite{stefanucci_excitonic_2025}, we infer that
$N_{\l\blQ}(\t)\propto \d_{\l\l_{0}}\d_{\blQ\bz}$ is  independent of $\t$.
Thus, in this scenario, the TR-ARPES signal is identical
in the coherent and incoherent
regimes~\cite{notepolaronic}.
Revealing the coherence of the system with TR-ARPES
requires a
probe field with a duration comparable to the exciton
period~\cite{perfetto_time-resolved_2020}.

An important feature of Eq.~(\ref{G<tmatfinld}) is that the total
number of carriers in the conduction band is given by
\begin{align}
N_{c}(\t)=-i\int\frac{d\w}{2\p}\sum_{c\blk}G^{<}_{cc\blk}(\t,\w)=
\sum_{\l\blQ}N_{\l\blQ}(\t),
\label{nc=nx}
\end{align}
where the normalization of the exciton wavefunctions
$\sum_{cv\blk}|A^{\l\blQ}_{cv\blk}|^{2}=1$ has been used.
In other words, in the low density limit $N_{c}$ is the same as the
total number of excitons, as it should be.
We remark that this property would not be satisfied if we  had
included only the T-matrix diagram in the self-energy.

{\em Beyond low density.--}
The formula Eq.~(\ref{G<}) for $G^{<}$  simplifies if the wavefunction
of the occupied exciton states
have a predominant contribution on a single conduction
band and valence band, say $c_{0}$ and $v_{0}$. In this case,
$K^{\l\blQ}_{cv\blk}=\d_{cc_{0}}\d_{vv_{0}}K^{\l\blQ}_{c_{0}v_{0}\blk}$,
and the quantity $\s$ in Eq.~(\ref{sigmasmall}) has one single
nonvanishing entry for $c=c'=c_{0}$. This implies that also the
self-energy is nonvanishing only for $c=c'=c_{0}$. Defining
$\D_{\blk}^{\l\blQ}(\t)\equiv
\sqrt{N_{\l\blQ}(\t)}K_{c_{0}v_{0}\blk}^{\l\blQ}$, and omitting the
depence of the various quantities on $\t$, we have
\begin{align}
G^{R}_{c_{0}c_{0}\blk}(\w)=\frac{1}{\w-\e_{c_{0}\blk}-
\sum_{\l\blQ}
\frac{|\D_{\blk}^{\l\blQ}|^{2}}{\w-E_{\l\blQ}-\e_{v_{0}\blk-\blQ}+i\eta}+i\eta}.
\label{twoband}
\end{align}
To gain insight into this problem,  we assume that
only a discrete number of excitons $\l_{j}\blQ_{j}$ are populated,
and that $ \D_{\blk}^{j}\equiv \D_{\blk}^{\l_{j}\blQ_{j}}$ are nonvanishing in small and
nonoverlapping regions $\callR^{j}$
of the first Brillouin zone. Then, for  $\blk\in \callR^{j}$ only
$\l_{j}\blQ_{j}$ contributes to the sum in the denominator of
Eq.~(\ref{twoband}). This allows us to calculate
$G^{<}_{c_{0}c_{0}\blk}$ analytically, see Appendix~\ref{hsextwobandapp} for details, and
find
\begin{align}
G^{<}_{c_{0}c_{0}\blk}(\w)\simeq 2\p i  \sum_{j}
\frac{E_{\l_{j}\blQ_{j}}+\e_{v_{0}\blk-\blQ_{j}}-E_{\blk,j}^{-}}
{E_{\blk,j}^{+}-E_{\blk,j}^{-}}\d(\w-E_{\blk,j}^{-}),
\label{G<i2}
\end{align}
where
\begin{align}
E_{\blk,j}^{\pm}&=\frac{1}{2}
\Big[\e_{c_{0}\blk}+E_{\l_{j}\blQ_{j}}+\e_{v_{0}\blk-\blQ_{j}}
\nn\\
&\pm\sqrt{(\e_{c_{0}\blk}-E_{\l_{j}\blQ_{j}}
-\e_{v_{0}\blk-\blQ_{j}})^{2}+4|\D_{\blk}^{j}|^{2}}\Big].
\label{epm}
\end{align}


To lowest order in $|\D^{j}_{\blk}|$
(low density limit) we can approximate the argument of the
Dirac-delta as $E_{\blk,j}^{-}\simeq
E^{\rm eq}_{\l_{j}\blQ_{j}}+\e^{\rm eq}_{v_{0}\blk-\blQ_{j}}$,
and the prefactor of the Dirac-delta as
$|\D_{\blk}^{j}|^{2}/(\e^{\rm eq}_{c_{0}\blk}-E^{\rm eq}_{\l_{j}\blQ_{j}}
-\e^{\rm eq}_{v_{0}\blk-\blQ_{j}})^{2}=N_{\l_{j}\blQ_{j}}|A^{\l_{j}\blQ_{j}}_{c_{0}v_{0}\blk}|^{2}$;
the resulting $G^{<}$ correctly agrees with
Eq.~(\ref{G<tmatfinld}).
Increasing the excitation density, the nonequilibrium
energy $E_{\l_{j}\blQ_{j}}(\t)+\e_{v_{0}\blk-\blQ_{j}}(\t)$
might overcome $\e_{c_{0}\blk}(\t)$, leading to a
hybridization (proportional to $|\D^{j}_{\blk}|$) between the
conduction band and the replica of the valence band, see
Fig.~\ref{spectrafig}(f).
This scenario has been thoroughly investigated in the
coherent
regime~\cite{perfetto_pump-driven_2019,perfetto_time-resolved_2020,chan_giant_2023} -- hence for $\blQ=\bz$ --
and is associated with a BEC-BCS crossover of the coherent exciton
fluid~\cite{ostreich_non-stationary_1993,hannewald_excitonic_2000}.
Our result in Eq.~(\ref{G<i2})
indicates that the same phenomenon can also occur in the
incoherent regime. In particular, the excitonic sidebands induced by
a sizable population of incoherent
excitons of momentum $\blQ$ might resemble the conduction
bands hybridized with a $\blQ$-shifted replica of the valence bands.



{\em TR-ARPES in WSe$_{2}$ monolayer.--} We consider a WSe$_{2}$
monolayer (ML) at low ($n_{c}\sim 10^{11}$~cm$^{-2}$) and moderate
($n_{c}\sim 10^{13}$~cm$^{-2}$) excitation densities $n_{c}=N_{c}/A$,
where $A$ is the area of the sample. The WSe$_{2}$ ML is a direct gap semiconductor,
with a
calculated bandgap at the $K$ and $K'$ points [Fig.~\ref{spectrafig}(a)]
estimated to be $\sim 2$~eV~\cite{lin_narrow-band_2021}.
For sufficiently long times after pumping, the exciton populations follow a
Bose distribution
$N_{\l\blQ}=[e^{(E_{\l\blQ}-\m_{X})/T}-1]^{-1}$, where $T$ is the
temperature and $\m_{\rm X}$ is self-consistently calculated from
Eq.~(\ref{nc=nx}).
In Fig.~\ref{spectrafig}(b) we show the equilibrium
exciton band structure, and how excitons are distributed
at temperature $T=70$~K; numerical details are reported in
Appendix~\ref{numapp}.
The low-energy landascape is characterized by several (almost) degenerate
excitons~\cite{deilmann_finite-momentum_2019,madeo_directly_2020},
the lowest excitons being at $K$ and $K'$.

At zero temperature, only the degenerate excitons with momenta $\blQ$
at $K$ and $K'$ are populated. Along the path $\G-K-\G$, and
for low values of $n_{c}$, the spectrum features a single sideband
for momenta $\blk$ near $K$ [Fig.~\ref{spectrafig}(d)].
This is the replica of $\e_{v\blk-\blQ}$ with $\blQ$
at the inequivalent $K$ point, which is the same as
$\e_{v\blk}$ with $\blk$ around $K'$. The replica is shifted in
energy by $\sim E_{\l\blQ}^{\rm eq}$, and the
intensity of the signal is proportional to $|A^{\l\blQ}_{cv\blk-\blQ}|^{2}$, in
agreement with Eq.~(\ref{G<tmatfinld}). Increasing $n_{c}$,
the exciton energies undergo a blue shift while the
bandgap shrinks [Figs.~\ref{spectrafig}(c) and (f)].
The replica of the valence band hybridize with the
conduction band, and the excitonic sideband acquires the shape of a
mexican hat [well described by Eq.~(\ref{epm})], as shown in
Fig.~\ref{spectrafig}(g).

At temperatures $T=70$~K, exciton states with momenta near other high
symmetry points ($M$,
$Q$ and $\G$) also become
populated, see  Fig.~\ref{spectrafig}(b).
Along the path $\G-K-\G$, and
for low values of $n_{c}$, a second
replica of the valence band appears for momenta $\blk$ near
$Q$ [Fig.~\ref{spectrafig}(e)]. This emerges from
excitons with $-\blQ$ around $Q$, leading to a replica of $\e_{v\blk}$
with $\blk$ around $K$. As the exciton energies are almost
degenerate, the two replicas in
Fig.~\ref{spectrafig}(e) are located in the same spectral range.
Compared to Fig.~\ref{spectrafig}(d), the signal appears
blurred -- this is due to the presence of exciton states with momenta close to,
but not exactly at, the high symmetry points. The additional excitonic
sideband around the $Q$ point persists at higher excitation
densities, but its shape turns into a mexican hat, see
Fig.~\ref{spectrafig}(h), again in agreement with Eq.~(\ref{epm}).


{\em Conclusions.--} We presented a first-principles approach to
excitons in TR-ARPES, bridging the coherent and incoherent regimes
through a unified formula that depends solely on the total number of
excitons at the pump-probe delay $\t$.
Applicable across a wide
range of temperatures and
excitation densities, our formula provides a powerful
tool for understanding and engineering nonequilibrium band structures
of excitonic materials.
Band hybridization with dark excitons enables
controlled shifts of the valence bands in both energy and momentum,
and to reshape them into a mexican-hat dispersion.
In particular,
the mechanism of
Ref.~\cite{foster_quench-induced_2014,perfetto_floquet_2020}
naturally extends to  the incoherent regime, enabling
control over the topological character of excited materials
through exciton wavefunction symmetries.
This points to a
promising avenue for realizing nonequilibrium
topological phases with
lifetimes reaching the picosecond timescale.

\acknowledgments

We acknowledge insightful discussions with Keshav Dani and Felipe da
Jornada.
We acknowledge funding from Ministero Universit\`a e
Ricerca PRIN under grant agreement No. 2022WZ8LME,
from INFN through project TIME2QUEST,
from European Research Council MSCA-ITN TIMES under grant agreement 101118915,
and from Tor Vergata University through project TESLA.



\appendix
\begin{widetext}

\section{Finite excitation-density corrections}
\label{findensapp}

We here describe a practical method to estimate the finite-density
corrections to the quasi-particle and exciton energies from the sole
knowledge of the exciton populations $N_{\l\blQ}$.

Let $\r_{\m\m'\blk}(t)=-iG_{\m\m'\blk}^{<}(t,t)$ be the one-electron
density matrix at time $t$. We define
$\r^{\rm eq}_{\m\m'\blk}=\d_{\m\m'}f^{\rm eq}_{\m\blk}$ as the
equilibrium one-electron  density matrix, with equilibrium occupations $f^{\rm
eq}_{v\blk}=1$ and $f^{\rm eq}_{c\blk}=0$. The change of the density
matrix  caused by an external pump pulse is then
\begin{align}
\D\r_{\m\m'\blk}(t)=\r_{\m\m'\blk}(t)-\r^{\rm eq}_{\m\m'\blk}.
\end{align}
This change induces the following change in the HSEX potential
\begin{subequations}
\begin{align}
V_{cc'\blk}(t)&=\sum_{vv'\blk'}\big(
v_{c\blk\, v\blk'\, v'\blk' \,c'\blk}-W_{c\blk\, v\blk'\,  c'\blk\, v'\blk'}
\big)\D\r_{v'v\blk'}(t)
+\sum_{c_{1}c_{2}\blk'}\big(
v_{c\blk\, c_{1}\blk'\, c_{2}\blk' \,c'\blk}-W_{c\blk\, c_{1}\blk'\,
c'\blk\, c_{2}\blk'}
\big)\D\r_{c_{2}c_{1}\blk'}(t),
\\
V_{vv'\blk}(t)&=\sum_{cc'\blk'}\big(
v_{v\blk\, c\blk'\, c'\blk'\, v'\blk}-W_{v\blk\, c\blk'\,  v'\blk\, c'\blk'}
\big)\D\r_{c'c\blk'}(t)
+\!\sum_{v_{1}v_{2}\blk'}\!\big(
v_{v\blk\, v_{1}\blk'\, v_{2}\blk'\, v'\blk}-W_{v\blk\, v_{1}\blk'\,
v'\blk\, v_{2}\blk'}
\big)\D\r_{v_{2}v_{1}\blk'}(t),
\\
V_{cv\blk}(t)&=\sum_{v'c'\blk'}\big(
v_{c\blk\, v'\blk'\, c'\blk'\, v\blk}-W_{c\blk\, v'\blk'\,  v\blk\, c'\blk'}
\big)\D\r_{c'v'\blk'}(t)
=-\sum_{v'c'\blk'}K^{\bz}_{cv\blk,c'v'\blk'}\D\r_{c'v'\blk'}(t)
=V_{vc\blk}^{\ast}(t).
\label{hsexpotcv}
\end{align}
\label{hsexpot}
\end{subequations}
In the HSEX approximation the renormalization of the quasiparticle
energies can be estimated as
\begin{align}
\e_{\m\blk}=\e^{\rm eq}_{\m\blk}+V_{\m\m\blk},
\label{rqpene}
\end{align}
where
\begin{subequations}
\begin{align}
V_{cc\blk}&\simeq \sum_{v'\blk'}\big(
v_{c\blk\, v'\blk'\, v'\blk' \,c\blk}-W_{c\blk\, v'\blk'\,  c\blk\, v'\blk'}
\big)(f_{v'\blk'}-1)+\sum_{c'\blk'}\big(
v_{c\blk\, c'\blk'\, c'\blk' \,c\blk}-W_{c\blk\, c'\blk'\,
c\blk\, c'\blk'}
\big)f_{c'\blk'},
\label{epscR}
\\
V_{vv\blk}&\simeq \sum_{c'\blk'}\big(
v_{v\blk\, c'\blk'\, c'\blk'\, v\blk}-W_{v\blk\, c'\blk'\,  v\blk\, c'\blk'}
\big)f_{c'\blk'}
+\sum_{v'\blk'}\big(
v_{v\blk\, v'\blk'\, v'\blk'\, v\blk}-W_{v\blk\, v'\blk'\,
v\blk\, v'\blk'}\big)(f_{v'\blk'}-1).
\label{epsvR}
\end{align}
\label{epscvR}
\end{subequations}
In Eqs.~(\ref{epscvR}), the off-diagonal contributions of the density
matrix are discarded, i.e.,
$\D\r_{cc'\blk}\simeq \d_{cc'}f_{c\blk}$ and
$\D \r_{vv'\blk}\simeq \d_{vv'}(f_{v\blk}-1)$.
This is consistent with ignoring
finite-density corrections to the
quasiparticle wavefunctions. Thus, Eq.~(\ref{epscvR}) provides a
formula for the renormalized quasi-particle
energies in terms of the nonequilibrium occupations $f_{\m\blk}$.

Let us now come to the renormalization of the exciton energies.
At finite excitation density, the BSE for the
exciton energies and wavefunctions reads
\begin{align}
\sum_{c'v'\blk'}H^{\blQ}_{cv\blk,c'v'\blk'}
\frac{A^{\l\blQ}_{c'v'\blk'}}
{\sqrt{f_{v'\blk'}-f_{c'\blk'+\blQ}}}=
E_{\l\blQ}\frac{A^{\l\blQ}_{cv\blk}}{\sqrt{f_{v\blk}-f_{c\blk+\blQ}}},
\label{bse1}
\end{align}
with
\begin{align}
H^{\blQ}_{cv\blk,c'v'\blk'}&=\d_{cv\blk,c'v'\blk'}
(\e_{c\blk+\blQ}-\e_{v\blk})
-\sqrt{f_{v\blk}-f_{c\blk+\blQ}}\,K^{\blQ}_{cv\blk,c'v'\blk'}
\sqrt{f_{v'\blk'}-f_{c'\blk'+\blQ}}
\label{bse2}
\end{align}
the BSE Hamiltonian.
We rewrite Eq.~(\ref{bse2}) as
\begin{align}
H^{\blQ}_{cv\blk,c'v'\blk'}=H^{{\rm eq},\blQ}_{cv\blk,c'v'\blk'}+\D
H^{\blQ}_{cv\blk,c'v'\blk'},
\end{align}
where $H^{{\rm eq},\blQ}$ is the equilibrium BSE Hamiltonian,
calculated with equilibrium occupations $f^{\rm eq}_{v\blk}=1$ and
$f^{\rm eq}_{c\blk}=0$, and
equilibrium band structure $\e_{\m\blk}=\e^{\rm eq}_{\m\blk}$.
To lowest order in $\D H^{\blQ}$ we have
\begin{align}
E_{\l\blQ}=E^{\rm eq}_{\l\blQ}+\sum_{cv\blk}\sum_{c'v'\blk'}
A^{\l\blQ\ast}_{cv\blk}\D
H^{\blQ}_{cv\blk,c'v'\blk'}A^{\l\blQ}_{c'v'\blk'},
\label{rxene}
\end{align}
where $E^{\rm eq}_{\l\blQ}$ is the equilibrium exciton energy and
$A^{\l\blQ}_{cv\blk}$ is the unperturbed exciton wavefunction.

We express $\D H^{\blQ}$ in terms of the nonequilibrium occupations
$f_{\m\blk}$.
For the Pauli blocking factors we have
\begin{align}
\sqrt{f_{v\blk}-f_{c\blk+\blQ}}
\simeq
1+\frac{1}{2}\big(f_{v\blk}-f_{c\blk+\blQ}-1\big).
\end{align}
We conclude that to lowest order in the excitation density
\begin{align}
\D H^{\blQ}_{cv\blk,c'v'\blk'}&=
\d_{cv\blk,c'v'\blk'}\Big[V_{cc\blk+\blQ}-V_{vv\blk}\Big]
-\frac{1}{2}\big(f_{v\blk}-f_{c\blk+\blQ}+
f_{v\blk'}-f_{c\blk'+\blQ}-2\big)
K^{\blQ}_{cv\blk,c'v'\blk'},
\label{deltabse}
\end{align}
where $V_{\m\m\blk}$ is given in Eq.~(\ref{epscvR}).

Equations~(\ref{rqpene}) and (\ref{rxene}) provide an approximation for the
nonequilibrium quasi-particle energies and exciton energies in terms
of the
nonequilibrium occupations. To  express these occupations in terms of exciton
populations we observe that for small excitation densities
the contribution of the $\l\blQ$-exciton to  $f_{c\blk}$ is
\begin{align}
\bra\l\blQ|\hat{d}^{\dag}_{c\blk}
\hat{d}_{c\blk}|\l\blQ\ket
&=\sum_{c_{1}v_{1}\blk_{1}}\sum_{c_{2}v_{2}\blk_{2}}
A^{\l\blQ\ast}_{c_{1}v_{1}\blk_{1}}
A^{\l\blQ}_{c_{2}v_{2}\blk_{2}}
\bra\Q_{g}|\hat{d}^{\dag}_{v_{1}\blk_{1}}
\hat{d}_{c_{1}\blk_{1}+\blQ}\hat{d}^{\dag}_{c\blk}
\hat{d}_{c\blk}\hat{d}^{\dag}_{c_{2}\blk_{2}+\blQ}
\hat{d}_{v_{2}\blk_{2}}|\Q_{g}\ket
=\sum_{v}|A^{\l\blQ}_{cv\blk-\blQ}|^{2}.
\end{align}
Therefore
\begin{align}
f_{c\blk}=\sum_{\l\blQ
v}N_{\l\blQ}|A^{\l\blQ}_{cv\blk-\blQ}|^{2},
\label{fck}
\end{align}
and similarly
\begin{align}
f_{v\blk}=1-\sum_{\l\blQ
c}N_{\l\blQ}|A^{\l\blQ}_{cv\blk}|^{2}.
\label{fvk}
\end{align}
Notice that the normalization of the exciton wavefunctions
guarantees that
$\sum_{c\blk}f_{c\blk}=\sum_{\l\blQ}N_{\l\blQ}$.

In Fig.~\ref{finite-corrections} we compare the equilibrium band structure and exciton
dispersions with those at moderate excitation density $n_{c}\simeq
10^{13}$~cm$^{-2}$ for the monolayer WSe$_{2}$. Due to finite density corrections
the replica of high-energy valence bands intersect the low-energy
conduction bands. In this scenario, the exciton-induced hybridization
causes the typical mexican hat shape in TR-ARPES.

\begin{figure}[t]
    \centering
\includegraphics[width=0.7\textwidth]{finite-corrections.pdf}
\caption{Equilibrium and nonequilibrium [excitation density $n\simeq
10^{13}$~cm$^{-2}$] exciton dispersions (a)
and quasiparticle energies (b). To highlight the band crossing between
the conduction bands and the replica of the valence bands, dashed
horizontal lines have been superimposed in correspondence of the
energy of the lowest exciton with vanishing momentum.}
\label{finite-corrections}
\end{figure}


\section{Lehmann approach}
\label{lehmannapp}
The state of the system just {\em after} pumping ($t>0$)
can be written as
\begin{align}
|\Q(t)\ket=|\Q_{g}\ket+\sum_{\l}\b_{\l}e^{-iE^{\rm eq}_{\l\bz}t}|\l\bz\ket.
\label{cohstate}
\end{align}
The coefficients $\b_{\l}$ are proportional to the excitonic
polarizations since
\begin{align}
\r_{\l}(t)&=
\sum_{cv\blk}A^{\l\bz\ast}_{cv\blk}\,
\bra\Q(t)|\hat{d}^{\dag}_{v\blk}\hat{d}_{c\blk}|\Q(t)\ket
=\sum_{cv\blk}A^{\l\bz\ast}_{cv\blk}\,
\sum_{\l'}\b_{\l'}A^{\l'\bz}_{cv\blk}e^{-iE^{\rm eq}_{\l'\bz}t}
=\b_{\l}e^{-iE^{\rm eq}_{\l\bz}t}.
\label{rldef}
\end{align}
Thus, Eq.~(\ref{cohstate}) can also be written as
\begin{align}
|\Q(t)\ket=|\Q_{g}\ket+\sum_{\l}\r_{\l}(t)|\l\bz\ket.
\label{cohstate2}
\end{align}
For a pure state the Green's function is given by
\begin{align}
G^{<}_{\m\m'\blk}(t,t')\equiv i\bra\Q(t')|\hat{d}^{\dag}_{\m'\blk}e^{i\hat{H}(t-t')}
\hat{d}_{\m\blk}|\Q(t)\ket,
\label{g<def}
\end{align}
where $\hat{H}$ is the many-body Hamiltonian of the unperturbed system.
We have
\begin{align}
\hat{d}_{c'\blk}|\Q(t)\ket=
\sum_{\l}\b_{\l} e^{-iE^{\rm eq}_{\l\bz}t}\sum_{v}
A^{\l\bz}_{c'v\blk}\hat{d}_{v\blk}|\Q_{g}\ket.
\end{align}
Taking into account that $\hat{H}\hat{d}_{v\blk}|\Q_{g}\ket
=-\e^{\rm eq}_{v\blk}\hat{d}_{v\blk}|\Q_{g}\ket$,
we find for the Green's function in Eq.~(\ref{g<def})
\begin{align}
G^{<}_{cc'\blk}(t,t')=i\sum_{\l\l'v}\b_{\l}^{\ast}\b_{\l'}\;
A^{\l\bz\ast}_{c'v\blk}A^{\l'\bz}_{cv\blk}\;
e^{iE^{\rm eq}_{\l\bz}t'}e^{-iE^{\rm eq}_{\l'\bz}t}
e^{-i\e^{\rm eq}_{v\blk}(t-t')}.
\label{gfcohleh}
\end{align}
This Green's function depends on $t$ and $t'$ separately, due to the
interference of different exciton
states~\cite{rustagi_coherent_2019,zhang_long-lived_2025}. The Fourier
transform with respect to $t-t'$ is dominated by the diagonal
contributions and reads
\begin{align}
G^{<}_{cc'\blk}(\w)=2\p i\sum_{\l v}|\r_{\l}|^{2}\;
A^{\l\bz\ast}_{c'v\blk}A^{\l\bz}_{cv\blk}\;
\d(\w-E^{\rm eq}_{\l\bz}-\e^{\rm eq}_{v\blk}).
\label{gfcohlehw}
\end{align}

For low excitation densities, the coherence between $|\Q_{g}\ket$ and the exciton states
$|\l\bz\ket$ is suppressed by electron-phonon scatterings on a
timescale that varies from a few tens to a few hundreds of
femtoseconds (depending on temperature and excitation
densities)~\cite{moody_intrinsic_2015,selig_exciton_2016,jakubczyk_impact_2018}. In the incoherent regime the
polarization $\r_{\l}$ vanishes, and the
electronic subsystem is described by an ensemble of bright and dark
incoherent excitons. From the many-body density matrix of the ensemble
$\hat{\r}=|\Q_{g}\ket\bra\Q_{g}|+\sum_{\l\blQ}N^{\rm
inc}_{\l\blQ}|\l\blQ\ket\bra\l\blQ|$, with $N^{\rm inc}_{\l\blQ}$
the number of incoherent excitons,  we can calculate
the Green's function  according to
\begin{align}
G^{<}_{cc'\blk}(t,t')=i\Tr\big[\hat{\r}\,e^{i\hat{H}t'}\hat{d}^{\dag}_{\m'\blk}e^{i\hat{H}(t-t')}
\hat{d}_{\m\blk}e^{-i\hat{H}t}\big]=
\sum_{\l\blQ}N^{\rm
inc}_{\l\blQ}G^{\l\blQ<}_{cc'\blk}(t,t'),
\label{g<ens}
\end{align}
where
\begin{align}
G^{\l\blQ<}_{cc'\blk}(t,t')&=i\bra\l\blQ(t')|
\hat{d}^{\dag}_{c'\blk}e^{i\hat{H}(t-t')}\hat{d}_{c\blk}|\l\blQ(t)\ket
\nn\\
&=ie^{-iE^{\rm eq}_{\l\blQ}(t-t')}
\bra\l\blQ|\hat{d}^{\dag}_{c'\blk}e^{i\hat{H}(t-t')}\hat{d}_{c\blk}|\l\blQ\ket
\nn\\
&=ie^{-iE^{\rm eq}_{\l\blQ}(t-t')}\sum_{v}
A^{\l\blQ\ast}_{c'v\blk-\blQ}A^{\l\blQ}_{cv\blk-\blQ}
e^{-i\e^{\rm eq}_{v\blk-\blQ}(t-t')}.
\label{G<isht}
\end{align}
In the incoherent regime the Green's function depends exclusively on
the time difference, and its Fourier transform reads
\begin{align}
G^{<}_{cc'\blk}(\w)=2\p i\sum_{\l\blQ v}N^{\rm inc}_{\l\blQ}
A^{\l\blQ\ast}_{c'v\blk-\blQ}A^{\l\blQ}_{cv\blk-\blQ}\;
\d(\w-E^{\rm eq}_{\l\blQ}-\e^{\rm eq}_{v\blk-\blQ}).
\label{G<ishtw}
\end{align}
Notice that if the population of incoherent excitons is given by
$N^{\rm inc}_{\l\blQ}=\d_{\blQ,\bz}|\r_{\l}|^{2}$, then the Green's
function -- and, consequently, the TR-ARPES spectrum -- becomes
identical for both the coherent and incoherent regimes.
This can be verified by comparing Eq.~(\ref{G<ishtw}) with the expression for the
Green's function in the coherent regime,
i.e., Eq.~(\ref{gfcohlehw}). Equations~(\ref{gfcohlehw})
and~(\ref{G<ishtw}), derived from the Lehmann approach,
agree with the Green's function in the many-body diagrammatic
approximation HSEX+TX, see Eq.~(\ref{G<tmatfinld}) in the main text.

\section{HSEX self-energy for two bands}
\label{hsextwobandapp}

We consider a system with only one active conduction and valence
bands. The HSEX lesser and greater Green's functions for
resonant (not necessarily weak) pumping with the $\l$
exciton have been derived in Ref.~\cite{perfetto_pump-driven_2019}, and
for sufficiently low temperatures they have a simple analytic form
\begin{subequations}
\begin{align}
G^{<}_{cc\blk}(\w)&=2\p i
\frac{E_{\l\bz}+\e_{v\blk}-E^{-}_{\blk}}
{E^{+}_{\blk}-E^{-}_{\blk}}
\d(\w-E^{-}_{\blk}),
\label{G<ccprm}
\\
G^{>}_{cc\blk}(\w)&=2\p i \frac{E_{\l\bz}+\e_{v\blk}-E^{+}_{\blk}}
{E^{+}_{\blk}-E^{-}_{\blk}}
\d(\w-E^{+}_{\blk}),
\label{G>ccprm}
\end{align}
\label{Gccprm}
\end{subequations}
where
\begin{subequations}
\begin{align}
E^{\pm}_{\blk}&\equiv \frac{1}{2}\Big[
\e_{c\blk}+E_{\l\bz}+\e_{v\blk}
\pm S_{\blk}\Big],
\label{Eplusminus}
\\
S_{\blk}&\equiv\sqrt{(\e_{c\blk}-E_{\l\bz}-\e_{v\blk})^{2}
+4|\D_{\blk}|^{2}},
\\
\D_{\blk}&\equiv K^{\l\bz}_{cv\blk}\r_{\l}.
\end{align}
\end{subequations}
Setting $z=\w+i\eta$, the retarded Green's function is then given by
\begin{align}
G^{R}_{cc\blk}(\w)&=i\int\frac{d\w'}{2\p}
\frac{G^{>}_{cc\blk}(\w')-G^{<}_{cc\blk}(\w')}{z-\w'}
=\frac{z-E_{\l\bz}-\e_{v\blk}}{
(z-E^{+}_{\blk})(z-E^{-}_{\blk})}
=\frac{1}{z-\e_{c\blk}-\frac{|\D_{\blk}|^{2}}{z-E_{\l\bz}-\e_{v\blk}}},
\label{GRanalyt2}
\end{align}
from which we can read the retarded HSEX self-energy
\begin{align}
\S^{R}_{cc\blk}(\w)=\frac{|\D_{\blk}|^{2}}{z-E_{\l\bz}-\e_{v\blk}}-i\eta.
\label{twobself}
\end{align}

To obtain the lesser HSEX self-energy we observe that
the Green's function in Eq.~(\ref{G<ccprm}) can also be
written as
\begin{align}
G^{<}_{cc\blk}(\w)=G^{R}_{cc\blk}(\w)\S^{<}_{cc\blk}(\w)G^{A}_{cc\blk}(\w).
\label{G<ccguess}
\end{align}
Let us show that this relation is satisfied for
\begin{align}
\S^{<}_{cc\blk}(\w)=-f(\w-E_{\l\bz}-\e_{v\blk})
\big[\S^{R}_{cc\blk}(\w)-\S^{A}_{cc\blk}(\w)\big],
\label{twobself2}
\end{align}
with $f(\w)=1/[e^{\w/T}+1]$ the Fermi function.
We have
\begin{align}
G^{R}_{cc\blk}(\w)\big[\S^{R}_{cc\blk}(\w)-\S^{A}_{cc\blk}(\w)\big]
G^{A}_{cc\blk}(\w)=G^{R}_{cc\blk}(\w)-G^{A}_{cc\blk}(\w)=
G^{>}_{cc\blk}(\w)-G^{<}_{cc\blk}(\w).
\end{align}
Therefore
\begin{align}
G^{R}_{cc\blk}(\w)\S^{<}_{cc\blk}(\w)G^{A}_{cc\blk}(\w)=
-2\p i f(\w-E_{\l\bz}-\e_{v\blk})
\Big[\frac{E_{\l\bz}+\e_{v\blk}-E^{+}_{\blk}}
{E^{+}_{\blk}-E^{-}_{\blk}}
\d(\w-E^{+}_{\blk})-\frac{E_{\l\bz}+\e_{v\blk}-E^{-}_{\blk}}
{E^{+}_{\blk}-E^{-}_{\blk}}
\d(\w-E^{-}_{\blk})\Big].
\label{gccguess2}
\end{align}
Next we observe that Eq.~(\ref{Eplusminus}) implies
\begin{align}
E^{\pm}_{\blk}-E_{\l\bz}-\e_{v\blk}&=
\frac{1}{2}\Big[
\e_{c\blk}+E_{\l\bz}+\e_{v\blk}
\pm S_{\blk}\Big]-E_{\l\bz}-\e_{v\blk}
\nn\\
&=\frac{1}{2}\Big[
\e_{c\blk}-E_{\l\bz}-\e_{v\blk}
\pm \sqrt{(\e_{c\blk}-E_{\l\bz}-\e_{v\blk})^{2}
+4\D_{\blk}^{2}}\Big]\gtrless 0
\end{align}
We conclude the second term in Eq.~(\ref{gccguess2}) dominates, and
we recover the analytic result in Eq.~(\ref{G<ccprm}).

\section{HSEX self-energy for multiple bands}
\label{hsexmultbandapp}

In the general case, the Green's function in the HSEX approximation
can be obtained from the
equation of motion on the Keldysh contour~\cite{svl-book_2025}
\begin{align}
\Big[i\frac{d}{dz}-\e^{\rm eq}_{\m\blk}\Big]G_{\m\m'\blk}(z,z')
-\sum_{\n}V_{\m\n\blk}(t)G_{\n\m'\blk}(z,z')=\d(z,z'),
\label{eomGRmatk}
\end{align}
where the HSEX potential is given in Eqs.~(\ref{hsexpot}).
We look for an approximate solution of Eq.~(\ref{eomGRmatk}) in
the conduction subspace, which is the subspace we are interested in.
We discard
all matrix elements $V_{cc'}$ with $c\neq c'$
and $V_{vv'}$ with $v\neq v'$. This is equivalent to discard
finite-density corrections to the
quasiparticle wavefunctions. Although these corrections do not pose a conceptual
problem, they considerably complicate the practical implementation of
the theory. From Eq.~(\ref{eomGRmatk}) we  have
\begin{subequations}
\begin{align}
\Big[i\frac{d}{dz}-\e_{c\blk}\Big]G_{cc'\blk}(z,z')
-\sum_{v}V_{cv\blk}(t)G_{vc'\blk}(z,z')&=\d(z,z')	,
\label{eomGccpk}
\end{align}
\begin{align}
\Big[i\frac{d}{dz}-\e_{v\blk}\Big]G_{vc'\blk}(z,z')
-\sum_{c''}V_{vc''\blk}(t)G_{c''c'\blk}(z,z')&=0	,
\end{align}
\end{subequations}
where  $\e_{\m\blk}\equiv \e^{\rm eq}_{\m\blk}+V_{\m\m\blk}$ are the
renormalized quasiparticle energies.
The second equation is solved by
\begin{align}
G_{vc'\blk}(z,z')=\int d\bar{z}\;g_{v\blk}(z,\bar{z})
\sum_{c''}V_{vc''\blk}(\bar{t})G_{c''c'\blk}(\bar{z},z'),
\label{solGvc}
\end{align}
where $g_{v\blk}$ is the solution of
\begin{align}
\Big[i\frac{d}{dz}-\e_{v\blk}\Big]g_{v\blk}(z,z')=\d(z,z')	.
\label{eomgvk}
\end{align}
Inserting Eq.~(\ref{solGvc}) into Eq.~(\ref{eomGccpk}) we find
\begin{align}
\Big[i\frac{d}{dz}-\e_{c\blk}\Big]G_{cc'\blk}(z,z')
-\sum_{c''}\int d\bar{z}\;\S_{cc''\blk}(z,\bar{z})
G_{c''c'\blk}(\bar{z},z')
=\d(z,z')	,
\label{eomGccp2}
\end{align}
where the HSEX self-energy reads
\begin{align}
\S_{cc'\blk}(z,z')&=
\sum_{v}V_{cv\blk}(t) g_{v\blk}(z,z')
V^{\ast}_{c'v\blk}(t').
\label{sigmahsex1}
\end{align}

Equation~(\ref{sigmahsex1}) can be further manipulated.
The HSEX dynamics of the density matrix is governed by the equation of motion
$i\frac{d}{dt}\r_{\blk}=[h^{\rm HSEX}_{\blk},\r_{\blk}]$ with
$h^{\rm
HSEX}_{\m\n\blk}=\d_{\m\n}\e^{\rm eq}_{\m\blk}+V_{\m\n\blk}$.
Therefore we have
\begin{align}
i\frac{d}{dt}\r_{cv\blk}&=(\e^{\rm eq}_{c\blk}-\e^{\rm eq}_{v\blk})\r_{cv\blk}+
\sum_{c'}V_{cc'\blk}\r_{c'v\blk}-\sum_{v'}\r_{cv'\blk}V_{v'v\blk}
+
\sum_{v'}V_{cv'\blk}\r_{v'v\blk}-\sum_{c'}\r_{cc'\blk}V_{c'v\blk}
\nn\\
&\simeq
(\e_{c\blk}-\e_{v\blk})\r_{cv\blk}+(f_{v\blk}-f_{c\blk})V_{cv\blk},
\label{eomrfdk}
\end{align}
where we implement the previously introduced approximations  $V_{cc'\blk}\simeq
\d_{cc'}V_{cc\blk}$, $V_{vv'\blk}\simeq
\d_{vv'}V_{vv\blk}$, as well as
$\r_{cc'\blk}\simeq \d_{cc'}f_{c\blk}$,
$\r_{vv'\blk}\simeq \d_{vv'}f_{v\blk}$, $f_{\m\blk}$ being the
occupation of the quasiparticle state $\m\blk$.
Taking into account that the off-diagonal
HSEX potential $V_{cv\blk}(t)$ is given by Eq.~(\ref{hsexpotcv}),
and comparing with the BSE in Eq.~(\ref{bse1}), we
see that
Eq.~(\ref{eomrfdk}) is solved by
\begin{align}
\r_{cv\blk}(t)=\sum_{\l}A^{\l\bz}_{cv\blk}\r_{\l}e^{-iE_{\l\bz}t}.
\end{align}
Ignoring the finite-density corrections to the
exciton wavefunctions, we have
\begin{align}
\sum_{c'v'\blk'}K^{\bz}_{cv\blk,c'v'\blk'}
A^{\l\bz}_{c'v'\blk'}=(\e^{\rm eq}_{c\blk}-\e^{\rm
eq}_{v\blk}-E_{\l\bz}^{\rm eq})A^{\l\bz}_{cv\blk}=K^{\l\bz}_{cv\blk},
\end{align}
and therefore
$V_{cv\blk}(t)=\sum_{\l}K^{\l\bz}_{cv\blk}\r_{\l}e^{-iE_{\l\bz}t}$.
Substituing this result into Eq.~(\ref{sigmahsex1}) we obtain
\begin{align}
\S_{cc'\blk}(z,z')&=
\sum_{v\l\l'}
K^{\l\bz}_{cv\blk}\r_{\l}e^{-iE_{\l\bz}t}
g_{v\blk}(z,z')e^{iE_{\l'\bz}t'}
K^{\l'\bz\ast}_{c'v\blk}\r^{\ast}_{\l'}.
\label{shsexk}
\end{align}

Let us extract the retarded component of the self-energy. From Eq.~(\ref{eomgvk}) we have
\begin{align}
g_{v\blk}^{R}(t,t')=-i\th(t-t')e^{-i\e_{v\blk}(t-t')}.
\end{align}
For sufficiently long probes, beating terms can be
ignored~\cite{rustagi_coherent_2019,zhang_long-lived_2025}.
Retaining only terms with $\l=\l'$, the retarded component of
Eq.~(\ref{shsexk}) in frequency space reads
\begin{align}
\S^{R}_{cc'\blk}(\w)=
\sum_{\l v}
K_{cv\blk}^{\l\bz}\frac{|\r_{\l}|^{2}}{z-E_{\l\bz}-\e_{v\blk}}
K^{\l\bz\ast}_{c'v\blk}-i\eta.
\label{shsexk2}
\end{align}
This formula generalizes Eq.~(\ref{twobself}) to multiple bands.
Accordingly, the generalization of Eq.~(\ref{twobself2}) for the lesser
self-energy is
\begin{align}
\S^{<}_{cc'\blk}(\w)&=-\sum_{\l v}
f(\w-E_{\l\bz}-\e_{v\blk})
\Big[
K_{cv\blk}^{\l\bz}\Big(\frac{|\r_{\l}|^{2}}{z-E_{\l\bz}-\e_{v\blk}}
-\frac{|\r_{\l}|^{2}}{z^{\ast}-E_{\l\bz}-\e_{v\blk}}\Big)
K^{\l\bz\ast}_{c'v\blk}
-2i\eta\Big].
\label{twobselfgen}
\end{align}
Equations~(\ref{shsexk2}) and (\ref{twobselfgen}) are the
self-energies of the main text for the coherent regime, i.e.,
Eqs.~(\ref{sigma}) with $\s$ evaluated in
$N_{\l\blQ}=\d_{\blQ\bz}|\r_{\l}|^{2}$.

\section{TX self-energy}
\label{txapp}

We expand the xc function in the excitonic
basis~\cite{perfetto_first-principles_2016,stefanucci_excitonic_2025}
\begin{align}
L^{\blQ\lessgtr}_{c_{1}v_{1}\blk_{1},c_{2}v_{2}\blk_{2}}(t,t')
=\sum_{\l}A^{\l\blQ}_{c_{1}v_{1}\blk_{1}}
N^{\lessgtr}_{\l\blQ}e^{-iE_{\l\blQ}(t-t')}
A^{\l\blQ\ast}_{c_{2}v_{2}\blk_{2}},
\end{align}
with $N^{<}_{\l\blQ}=N^{\rm inc}_{\l\blQ}$ and
$N^{>}_{\l\blQ}=1+N^{\rm inc}_{\l\blQ}$, and approximate the
Green's function at the quasiparticle level,
$G^{\lessgtr}_{v\blk}(t,t')=if^{\lessgtr}_{v\blk}e^{-i\e_{v\blk}(t-t')}$,
with $f^{<}_{v\blk}=f_{v\blk}$ and $f^{>}_{v\blk}=f_{v\blk}-1$.
Then, the Fourier transform of the lesser and greater TX self-energy,
see Eq.~(\ref{sigma<>}),
read
\begin{align}
\S^{\lessgtr}_{cc'\blk}(\w)&=2\p i
\sum_{\l\blQ v}K^{\l\blQ}_{cv\blk}\;N^{\lessgtr}_{\l\blQ}
f^{\lessgtr}_{v\blk-\blQ}\;
K^{\l\blQ\ast}_{c'v\blk}
 \d(\w-E_{\l\blQ}-\e_{v\blk-\blQ}).
\label{sigma3}
\end{align}
The retarded TX self-energy can be calculated from
\begin{align}
\S^{R}_{cc'\blk}(\w)&=i\int\frac{d\w'}{2\p}
\frac{\S^{>}_{cc'\blk}(\w)-\S^{<}_{cc'\blk}(\w)}{\w-\w'+i\eta}-i\eta
=\sum_{\l\blQ v}K^{\l\blQ}_{cv\blk}\;
\frac{N_{\l\blQ}+1-f_{v\blk-\blQ}}{\w-E_{\l\blQ}-\e_{v\blk-\blQ}+i\eta}\;
K^{\l\blQ\ast}_{c'v\blk}-i\eta.
\label{sigmaRapp}
\end{align}
For not too large excitation densities, we can approximate
$1-f_{v\blk-\blQ}\simeq 0$, thus recovering the result in the main
text.
In analogy with Eq.~(\ref{twobselfgen}), we approximate the lesser TX self-energy
as
\begin{align}
\S^{<}_{cc'\blk}(\w)=- \sum_{\l\blQ v}
f(\w-E_{\l\blQ}-\e_{v\blk-\blQ})
\Big[K^{\l\blQ}_{cv\blk}\Big(\frac{N_{\l\blQ}}{z-E_{\l\blQ}-\e_{v\blk-\blQ}}
-\frac{N_{\l\blQ}}{z^{\ast}-E_{\l\blQ}-\e_{v\blk-\blQ}}\Big)
K^{\l\blQ\ast}_{c'v\blk}-2i\eta\Big].
\label{twobselfgenic}
\end{align}
Equations~(\ref{sigmaRapp}) and (\ref{twobselfgenic}) are the
self-energies of the main text for the incoherent regime, i.e.,
Eqs.~(\ref{sigma}) with $\s$ evaluated in
$N_{\l\blQ}=N_{\l\blQ}^{\rm inc}$.

\section{Numerical details}
\label{numapp}
The electronic structure is described using a spin-resolved density functional theory (DFT) band dispersion
$\varepsilon_{n\mathbf{k}}$, mapped onto a tight-binding representation as detailed in Ref.~\cite{PhysRevB.88.085433}.
To match the experimentally measured quasiparticle band gap of about
$2.0~\mathrm{eV}$, a rigid scissor correction of
$0.5~\mathrm{eV}$ is applied to the conduction-band states $\varepsilon_{c\mathbf{k}}$
\cite{PhysRevLett.113.076802,Hsu_2019}.
The eigenstates of the Bloch Hamiltonian, denoted by $\mathbf{U}_{\mathbf{k}}$,
serve as the basis for constructing the Coulomb matrix elements that appear in both the
Bethe–Salpeter equation (BSE) kernel and the Hartree–screened-exchange (HSEX) Hamiltonian,
following the methodology of Ref.~\cite{perfetto_real-time_2023}.

The Coulomb interaction is modeled using the Rytova–Keldysh potential
\cite{keldysh,PhysRevB.84.085406}, assuming a dielectric constant $\varepsilon = 1$
and a screening length of
$r_0 = 45~\mathring{\rm A}$ \cite{stier_magnetooptics_2018}.
To treat the divergence of the Coulomb potential $v_{\mathbf{q}}$ near the $\Gamma$ point,
we adopt the regularization scheme introduced in Refs.~\cite{PhysRevB.88.245309,PhysRevB.97.205409},
which consists of averaging $v_{\mathbf{q}}$ over a small region $\Omega$ around $\mathbf{q}=0$.
The size of this region is determined by the spacing of the Brillouin-zone sampling grid.

In typical semiconductors, Coulomb matrix elements that involve processes changing the number of electrons
in the valence or conduction bands are vanishingly small and are therefore neglected
\cite{PhysRevB.93.205145}.
Since the Rytova–Keldysh interaction tensor $V$ in Eq.~(5) already incorporates
\emph{ab initio} static screening over the momentum scales relevant for excitonic physics,
the screened exchange interaction $W$ entering both the BSE kernel and the HSEX Hamiltonian
is approximated by $W \approx V$.
This simplification has been demonstrated to reproduce experimental exciton binding energies
with high accuracy in two-dimensional systems
\cite{steinhoff2,PhysRevLett.113.076802,PhysRevB.88.045318}.

In the calculations presented here, the electronic subspace is limited to the two uppermost valence bands
and the two lowest conduction bands.
The Brillouin zone is sampled using a uniform mesh comprising $3072$ $\mathbf{k}$ points.



\end{widetext}









\begin{thebibliography}{69}
\makeatletter
\providecommand \@ifxundefined [1]{
 \@ifx{#1\undefined}
}
\providecommand \@ifnum [1]{
 \ifnum #1\expandafter \@firstoftwo
 \else \expandafter \@secondoftwo
 \fi
}
\providecommand \@ifx [1]{
 \ifx #1\expandafter \@firstoftwo
 \else \expandafter \@secondoftwo
 \fi
}
\providecommand \natexlab [1]{#1}
\providecommand \enquote  [1]{``#1''}
\providecommand \bibnamefont  [1]{#1}
\providecommand \bibfnamefont [1]{#1}
\providecommand \citenamefont [1]{#1}
\providecommand \href@noop [0]{\@secondoftwo}
\providecommand \href [0]{\begingroup \@sanitize@url \@href}
\providecommand \@href[1]{\@@startlink{#1}\@@href}
\providecommand \@@href[1]{\endgroup#1\@@endlink}
\providecommand \@sanitize@url [0]{\catcode `\\12\catcode `\$12\catcode
  `\&12\catcode `\#12\catcode `\^12\catcode `\_12\catcode `\%12\relax}
\providecommand \@@startlink[1]{}
\providecommand \@@endlink[0]{}
\providecommand \url  [0]{\begingroup\@sanitize@url \@url }
\providecommand \@url [1]{\endgroup\@href {#1}{\urlprefix }}
\providecommand \urlprefix  [0]{URL }
\providecommand \Eprint [0]{\href }
\providecommand \doibase [0]{https://doi.org/}
\providecommand \selectlanguage [0]{\@gobble}
\providecommand \bibinfo  [0]{\@secondoftwo}
\providecommand \bibfield  [0]{\@secondoftwo}
\providecommand \translation [1]{[#1]}
\providecommand \BibitemOpen [0]{}
\providecommand \bibitemStop [0]{}
\providecommand \bibitemNoStop [0]{.\EOS\space}
\providecommand \EOS [0]{\spacefactor3000\relax}
\providecommand \BibitemShut  [1]{\csname bibitem#1\endcsname}
\let\auto@bib@innerbib\@empty

\bibitem [{\citenamefont {Onida}\ \emph {et~al.}(2002)\citenamefont {Onida},
  \citenamefont {Reining},\ and\ \citenamefont
  {Rubio}}]{onida_electronic_2002}
  \BibitemOpen
  \bibfield  {author} {\bibinfo {author} {\bibfnamefont {G.}~\bibnamefont
  {Onida}}, \bibinfo {author} {\bibfnamefont {L.}~\bibnamefont {Reining}},\
  and\ \bibinfo {author} {\bibfnamefont {A.}~\bibnamefont {Rubio}},\ }\bibfield
   {title} {\bibinfo {title} {Electronic excitations: density-functional versus
  many-body green's-function approaches},\ }\href
  {https://doi.org/10.1103/RevModPhys.74.601} {\bibfield  {journal} {\bibinfo
  {journal} {Rev. Mod. Phys.}\ }\textbf {\bibinfo {volume} {74}},\ \bibinfo
  {pages} {601} (\bibinfo {year} {2002})}\BibitemShut {NoStop}
\bibitem [{\citenamefont {Zhang}\ \emph {et~al.}(2022)\citenamefont {Zhang},
  \citenamefont {Pincelli}, \citenamefont {Jozwiak}, \citenamefont {Kondo},
  \citenamefont {Ernstorfer}, \citenamefont {Sato},\ and\ \citenamefont
  {Zhou}}]{zhang_angle-resolved_2022}
  \BibitemOpen
  \bibfield  {author} {\bibinfo {author} {\bibfnamefont {H.}~\bibnamefont
  {Zhang}}, \bibinfo {author} {\bibfnamefont {T.}~\bibnamefont {Pincelli}},
  \bibinfo {author} {\bibfnamefont {C.}~\bibnamefont {Jozwiak}}, \bibinfo
  {author} {\bibfnamefont {T.}~\bibnamefont {Kondo}}, \bibinfo {author}
  {\bibfnamefont {R.}~\bibnamefont {Ernstorfer}}, \bibinfo {author}
  {\bibfnamefont {T.}~\bibnamefont {Sato}},\ and\ \bibinfo {author}
  {\bibfnamefont {S.}~\bibnamefont {Zhou}},\ }\bibfield  {title} {\bibinfo
  {title} {Angle-resolved photoemission spectroscopy},\ }\href
  {https://doi.org/10.1038/s43586-022-00133-7} {\bibfield  {journal} {\bibinfo
  {journal} {Nature Reviews Methods Primers}\ }\textbf {\bibinfo {volume}
  {2}},\ \bibinfo {pages} {54} (\bibinfo {year} {2022})}\BibitemShut {NoStop}
\bibitem [{\citenamefont {Boschini}\ \emph {et~al.}(2024)\citenamefont
  {Boschini}, \citenamefont {Zonno},\ and\ \citenamefont
  {Damascelli}}]{boschini_time-resolved_2024}
  \BibitemOpen
  \bibfield  {author} {\bibinfo {author} {\bibfnamefont {F.}~\bibnamefont
  {Boschini}}, \bibinfo {author} {\bibfnamefont {M.}~\bibnamefont {Zonno}},\
  and\ \bibinfo {author} {\bibfnamefont {A.}~\bibnamefont {Damascelli}},\
  }\bibfield  {title} {\bibinfo {title} {Time-resolved arpes studies of quantum
  materials},\ }\href {https://doi.org/10.1103/RevModPhys.96.015003} {\bibfield
   {journal} {\bibinfo  {journal} {Rev. Mod. Phys.}\ }\textbf {\bibinfo
  {volume} {96}},\ \bibinfo {pages} {015003} (\bibinfo {year}
  {2024})}\BibitemShut {NoStop}
\bibitem [{\citenamefont {Xu}\ and\ \citenamefont
  {Zong}(2025)}]{xu_time-domain_2025}
  \BibitemOpen
  \bibfield  {author} {\bibinfo {author} {\bibfnamefont {C.}~\bibnamefont
  {Xu}}\ and\ \bibinfo {author} {\bibfnamefont {A.}~\bibnamefont {Zong}},\
  }\bibfield  {title} {\bibinfo {title} {Time-domain study of coupled
  collective excitations in quantum materials},\ }\href
  {https://doi.org/10.1038/s41535-025-00726-x} {\bibfield  {journal} {\bibinfo
  {journal} {npj Quantum Materials}\ }\textbf {\bibinfo {volume} {10}},\
  \bibinfo {pages} {21} (\bibinfo {year} {2025})}\BibitemShut {NoStop}
\bibitem [{\citenamefont {Madéo}\ \emph {et~al.}(2020)\citenamefont {Madéo},
  \citenamefont {Man}, \citenamefont {Sahoo}, \citenamefont {Campbell},
  \citenamefont {Pareek}, \citenamefont {Wong}, \citenamefont {Al-Mahboob},
  \citenamefont {Chan}, \citenamefont {Karmakar}, \citenamefont {Mariserla},
  \citenamefont {Li}, \citenamefont {Heinz}, \citenamefont {Cao},\ and\
  \citenamefont {Dani}}]{madeo_directly_2020}
  \BibitemOpen
  \bibfield  {author} {\bibinfo {author} {\bibfnamefont {J.}~\bibnamefont
  {Madéo}}, \bibinfo {author} {\bibfnamefont {M.~K.~L.}\ \bibnamefont {Man}},
  \bibinfo {author} {\bibfnamefont {C.}~\bibnamefont {Sahoo}}, \bibinfo
  {author} {\bibfnamefont {M.}~\bibnamefont {Campbell}}, \bibinfo {author}
  {\bibfnamefont {V.}~\bibnamefont {Pareek}}, \bibinfo {author} {\bibfnamefont
  {E.~L.}\ \bibnamefont {Wong}}, \bibinfo {author} {\bibfnamefont
  {A.}~\bibnamefont {Al-Mahboob}}, \bibinfo {author} {\bibfnamefont {N.~S.}\
  \bibnamefont {Chan}}, \bibinfo {author} {\bibfnamefont {A.}~\bibnamefont
  {Karmakar}}, \bibinfo {author} {\bibfnamefont {B.~M.~K.}\ \bibnamefont
  {Mariserla}}, \bibinfo {author} {\bibfnamefont {X.}~\bibnamefont {Li}},
  \bibinfo {author} {\bibfnamefont {T.~F.}\ \bibnamefont {Heinz}}, \bibinfo
  {author} {\bibfnamefont {T.}~\bibnamefont {Cao}},\ and\ \bibinfo {author}
  {\bibfnamefont {K.~M.}\ \bibnamefont {Dani}},\ }\bibfield  {title} {\bibinfo
  {title} {Directly visualizing the momentum-forbidden dark excitons and their
  dynamics in atomically thin semiconductors},\ }\href
  {https://doi.org/10.1126/science.aba1029} {\bibfield  {journal} {\bibinfo
  {journal} {Science}\ }\textbf {\bibinfo {volume} {370}},\ \bibinfo {pages}
  {1199} (\bibinfo {year} {2020})}\BibitemShut {NoStop}
\bibitem [{\citenamefont {Dendzik}\ \emph {et~al.}(2020)\citenamefont
  {Dendzik}, \citenamefont {Xian}, \citenamefont {Perfetto}, \citenamefont
  {Sangalli}, \citenamefont {Kutnyakhov}, \citenamefont {Dong}, \citenamefont
  {Beaulieu}, \citenamefont {Pincelli}, \citenamefont {Pressacco},
  \citenamefont {Curcio}, \citenamefont {Agustsson}, \citenamefont {Heber},
  \citenamefont {Hauer}, \citenamefont {Wurth}, \citenamefont {Brenner},
  \citenamefont {Acremann}, \citenamefont {Hofmann}, \citenamefont {Wolf},
  \citenamefont {Marini}, \citenamefont {Stefanucci}, \citenamefont {Rettig},\
  and\ \citenamefont {Ernstorfer}}]{dendzik_observation_2020}
  \BibitemOpen
  \bibfield  {author} {\bibinfo {author} {\bibfnamefont {M.}~\bibnamefont
  {Dendzik}}, \bibinfo {author} {\bibfnamefont {R.~P.}\ \bibnamefont {Xian}},
  \bibinfo {author} {\bibfnamefont {E.}~\bibnamefont {Perfetto}}, \bibinfo
  {author} {\bibfnamefont {D.}~\bibnamefont {Sangalli}}, \bibinfo {author}
  {\bibfnamefont {D.}~\bibnamefont {Kutnyakhov}}, \bibinfo {author}
  {\bibfnamefont {S.}~\bibnamefont {Dong}}, \bibinfo {author} {\bibfnamefont
  {S.}~\bibnamefont {Beaulieu}}, \bibinfo {author} {\bibfnamefont
  {T.}~\bibnamefont {Pincelli}}, \bibinfo {author} {\bibfnamefont
  {F.}~\bibnamefont {Pressacco}}, \bibinfo {author} {\bibfnamefont
  {D.}~\bibnamefont {Curcio}}, \bibinfo {author} {\bibfnamefont {S.~Y.}\
  \bibnamefont {Agustsson}}, \bibinfo {author} {\bibfnamefont {M.}~\bibnamefont
  {Heber}}, \bibinfo {author} {\bibfnamefont {J.}~\bibnamefont {Hauer}},
  \bibinfo {author} {\bibfnamefont {W.}~\bibnamefont {Wurth}}, \bibinfo
  {author} {\bibfnamefont {G.}~\bibnamefont {Brenner}}, \bibinfo {author}
  {\bibfnamefont {Y.}~\bibnamefont {Acremann}}, \bibinfo {author}
  {\bibfnamefont {P.}~\bibnamefont {Hofmann}}, \bibinfo {author} {\bibfnamefont
  {M.}~\bibnamefont {Wolf}}, \bibinfo {author} {\bibfnamefont {A.}~\bibnamefont
  {Marini}}, \bibinfo {author} {\bibfnamefont {G.}~\bibnamefont {Stefanucci}},
  \bibinfo {author} {\bibfnamefont {L.}~\bibnamefont {Rettig}},\ and\ \bibinfo
  {author} {\bibfnamefont {R.}~\bibnamefont {Ernstorfer}},\ }\bibfield  {title}
  {\bibinfo {title} {Observation of an excitonic mott transition through
  ultrafast core-cum-conduction photoemission spectroscopy},\ }\href
  {https://link.aps.org/doi/10.1103/PhysRevLett.125.096401} {\bibfield
  {journal} {\bibinfo  {journal} {Phys. Rev. Lett.}\ }\textbf {\bibinfo
  {volume} {125}},\ \bibinfo {pages} {096401} (\bibinfo {year}
  {2020})}\BibitemShut {NoStop}
\bibitem [{\citenamefont {Tanimura}\ and\ \citenamefont
  {Tanimura}(2020)}]{tanimura_time_2020}
  \BibitemOpen
  \bibfield  {author} {\bibinfo {author} {\bibfnamefont {H.}~\bibnamefont
  {Tanimura}}\ and\ \bibinfo {author} {\bibfnamefont {K.}~\bibnamefont
  {Tanimura}},\ }\bibfield  {title} {\bibinfo {title} {Time- and angle-resolved
  photoemission spectroscopy for the saddle-point excitons in {G}a{A}s},\
  }\href {https://doi.org/10.1103/PhysRevB.102.045204} {\bibfield  {journal}
  {\bibinfo  {journal} {Phys. Rev. B}\ }\textbf {\bibinfo {volume} {102}},\
  \bibinfo {pages} {045204} (\bibinfo {year} {2020})}\BibitemShut {NoStop}
\bibitem [{\citenamefont {Wallauer}\ \emph {et~al.}(2021)\citenamefont
  {Wallauer}, \citenamefont {Perea-Causin}, \citenamefont {Münster},
  \citenamefont {Zajusch}, \citenamefont {Brem}, \citenamefont {Güdde},
  \citenamefont {Tanimura}, \citenamefont {Lin}, \citenamefont {Huber},
  \citenamefont {Malic},\ and\ \citenamefont
  {Höfer}}]{wallauer_momentum-resolved_2021}
  \BibitemOpen
  \bibfield  {author} {\bibinfo {author} {\bibfnamefont {R.}~\bibnamefont
  {Wallauer}}, \bibinfo {author} {\bibfnamefont {R.}~\bibnamefont
  {Perea-Causin}}, \bibinfo {author} {\bibfnamefont {L.}~\bibnamefont
  {Münster}}, \bibinfo {author} {\bibfnamefont {S.}~\bibnamefont {Zajusch}},
  \bibinfo {author} {\bibfnamefont {S.}~\bibnamefont {Brem}}, \bibinfo {author}
  {\bibfnamefont {J.}~\bibnamefont {Güdde}}, \bibinfo {author} {\bibfnamefont
  {K.}~\bibnamefont {Tanimura}}, \bibinfo {author} {\bibfnamefont {K.-Q.}\
  \bibnamefont {Lin}}, \bibinfo {author} {\bibfnamefont {R.}~\bibnamefont
  {Huber}}, \bibinfo {author} {\bibfnamefont {E.}~\bibnamefont {Malic}},\ and\
  \bibinfo {author} {\bibfnamefont {U.}~\bibnamefont {Höfer}},\ }\bibfield
  {title} {\bibinfo {title} {Momentum-resolved observation of exciton formation
  dynamics in monolayer {WS}2},\ }\href
  {https://doi.org/10.1021/acs.nanolett.1c01839} {\bibfield  {journal}
  {\bibinfo  {journal} {Nano Letters}\ }\textbf {\bibinfo {volume} {21}},\
  \bibinfo {pages} {5867} (\bibinfo {year} {2021})}\BibitemShut {NoStop}
\bibitem [{\citenamefont {Fukutani}\ \emph {et~al.}(2021)\citenamefont
  {Fukutani}, \citenamefont {Stania}, \citenamefont {Il~Kwon}, \citenamefont
  {Kim}, \citenamefont {Kong}, \citenamefont {Kim},\ and\ \citenamefont
  {Yeom}}]{fukutani_detecting_2021}
  \BibitemOpen
  \bibfield  {author} {\bibinfo {author} {\bibfnamefont {K.}~\bibnamefont
  {Fukutani}}, \bibinfo {author} {\bibfnamefont {R.}~\bibnamefont {Stania}},
  \bibinfo {author} {\bibfnamefont {C.}~\bibnamefont {Il~Kwon}}, \bibinfo
  {author} {\bibfnamefont {J.~S.}\ \bibnamefont {Kim}}, \bibinfo {author}
  {\bibfnamefont {K.~J.}\ \bibnamefont {Kong}}, \bibinfo {author}
  {\bibfnamefont {J.}~\bibnamefont {Kim}},\ and\ \bibinfo {author}
  {\bibfnamefont {H.~W.}\ \bibnamefont {Yeom}},\ }\bibfield  {title} {\bibinfo
  {title} {Detecting photoelectrons from spontaneously formed excitons},\
  }\href {https://doi.org/10.1038/s41567-021-01289-x} {\bibfield  {journal}
  {\bibinfo  {journal} {Nature Physics}\ }\textbf {\bibinfo {volume} {17}},\
  \bibinfo {pages} {1024} (\bibinfo {year} {2021})}\BibitemShut {NoStop}
\bibitem [{\citenamefont {Kunin}\ \emph {et~al.}(2023)\citenamefont {Kunin},
  \citenamefont {Chernov}, \citenamefont {Bakalis}, \citenamefont {Li},
  \citenamefont {Cheng}, \citenamefont {Withers}, \citenamefont {White},
  \citenamefont {Sch\"onhense}, \citenamefont {Du}, \citenamefont {Kawakami},\
  and\ \citenamefont {Allison}}]{kunin_momentum-resolved_2023}
  \BibitemOpen
  \bibfield  {author} {\bibinfo {author} {\bibfnamefont {A.}~\bibnamefont
  {Kunin}}, \bibinfo {author} {\bibfnamefont {S.}~\bibnamefont {Chernov}},
  \bibinfo {author} {\bibfnamefont {J.}~\bibnamefont {Bakalis}}, \bibinfo
  {author} {\bibfnamefont {Z.}~\bibnamefont {Li}}, \bibinfo {author}
  {\bibfnamefont {S.}~\bibnamefont {Cheng}}, \bibinfo {author} {\bibfnamefont
  {Z.~H.}\ \bibnamefont {Withers}}, \bibinfo {author} {\bibfnamefont {M.~G.}\
  \bibnamefont {White}}, \bibinfo {author} {\bibfnamefont {G.}~\bibnamefont
  {Sch\"onhense}}, \bibinfo {author} {\bibfnamefont {X.}~\bibnamefont {Du}},
  \bibinfo {author} {\bibfnamefont {R.~K.}\ \bibnamefont {Kawakami}},\ and\
  \bibinfo {author} {\bibfnamefont {T.~K.}\ \bibnamefont {Allison}},\
  }\bibfield  {title} {\bibinfo {title} {Momentum-resolved exciton coupling and
  valley polarization dynamics in monolayer ${\mathrm{ws}}_{2}$},\ }\href
  {https://doi.org/10.1103/PhysRevLett.130.046202} {\bibfield  {journal}
  {\bibinfo  {journal} {Phys. Rev. Lett.}\ }\textbf {\bibinfo {volume} {130}},\
  \bibinfo {pages} {046202} (\bibinfo {year} {2023})}\BibitemShut {NoStop}
\bibitem [{\citenamefont {Bange}\ \emph {et~al.}(2024)\citenamefont {Bange},
  \citenamefont {Schmitt}, \citenamefont {Bennecke}, \citenamefont {Meneghini},
  \citenamefont {AlMutairi}, \citenamefont {Watanabe}, \citenamefont
  {Taniguchi}, \citenamefont {Steil}, \citenamefont {Steil}, \citenamefont
  {Weitz}, \citenamefont {Jansen}, \citenamefont {Hofmann}, \citenamefont
  {Brem}, \citenamefont {Malic}, \citenamefont {Reutzel},\ and\ \citenamefont
  {Mathias}}]{bange_probing_2024}
  \BibitemOpen
  \bibfield  {author} {\bibinfo {author} {\bibfnamefont {J.~P.}\ \bibnamefont
  {Bange}}, \bibinfo {author} {\bibfnamefont {D.}~\bibnamefont {Schmitt}},
  \bibinfo {author} {\bibfnamefont {W.}~\bibnamefont {Bennecke}}, \bibinfo
  {author} {\bibfnamefont {G.}~\bibnamefont {Meneghini}}, \bibinfo {author}
  {\bibfnamefont {A.}~\bibnamefont {AlMutairi}}, \bibinfo {author}
  {\bibfnamefont {K.}~\bibnamefont {Watanabe}}, \bibinfo {author}
  {\bibfnamefont {T.}~\bibnamefont {Taniguchi}}, \bibinfo {author}
  {\bibfnamefont {D.}~\bibnamefont {Steil}}, \bibinfo {author} {\bibfnamefont
  {S.}~\bibnamefont {Steil}}, \bibinfo {author} {\bibfnamefont {R.~T.}\
  \bibnamefont {Weitz}}, \bibinfo {author} {\bibfnamefont {G.~S.~M.}\
  \bibnamefont {Jansen}}, \bibinfo {author} {\bibfnamefont {S.}~\bibnamefont
  {Hofmann}}, \bibinfo {author} {\bibfnamefont {S.}~\bibnamefont {Brem}},
  \bibinfo {author} {\bibfnamefont {E.}~\bibnamefont {Malic}}, \bibinfo
  {author} {\bibfnamefont {M.}~\bibnamefont {Reutzel}},\ and\ \bibinfo {author}
  {\bibfnamefont {S.}~\bibnamefont {Mathias}},\ }\bibfield  {title} {\bibinfo
  {title} {Probing electron-hole coulomb correlations in the exciton landscape
  of a twisted semiconductor heterostructure},\ }\href
  {https://doi.org/10.1126/sciadv.adi1323} {\bibfield  {journal} {\bibinfo
  {journal} {Science Advances}\ }\textbf {\bibinfo {volume} {10}},\ \bibinfo
  {pages} {eadi1323} (\bibinfo {year} {2024})}\BibitemShut {NoStop}
\bibitem [{\citenamefont {Gosetti}\ \emph {et~al.}(2024)\citenamefont
  {Gosetti}, \citenamefont {Cervantes-Villanueva}, \citenamefont {Mor},
  \citenamefont {Sangalli}, \citenamefont {García-Cristóbal}, \citenamefont
  {Molina-Sánchez}, \citenamefont {Agekyan}, \citenamefont {Tuniz},
  \citenamefont {Puntel}, \citenamefont {Bronsch}, \citenamefont {Cilento},\
  and\ \citenamefont {Pagliara}}]{gosetti_unveiling_2024}
  \BibitemOpen
  \bibfield  {author} {\bibinfo {author} {\bibfnamefont {V.}~\bibnamefont
  {Gosetti}}, \bibinfo {author} {\bibfnamefont {J.}~\bibnamefont
  {Cervantes-Villanueva}}, \bibinfo {author} {\bibfnamefont {S.}~\bibnamefont
  {Mor}}, \bibinfo {author} {\bibfnamefont {D.}~\bibnamefont {Sangalli}},
  \bibinfo {author} {\bibfnamefont {A.}~\bibnamefont {García-Cristóbal}},
  \bibinfo {author} {\bibfnamefont {A.}~\bibnamefont {Molina-Sánchez}},
  \bibinfo {author} {\bibfnamefont {V.~F.}\ \bibnamefont {Agekyan}}, \bibinfo
  {author} {\bibfnamefont {M.}~\bibnamefont {Tuniz}}, \bibinfo {author}
  {\bibfnamefont {D.}~\bibnamefont {Puntel}}, \bibinfo {author} {\bibfnamefont
  {W.}~\bibnamefont {Bronsch}}, \bibinfo {author} {\bibfnamefont
  {F.}~\bibnamefont {Cilento}},\ and\ \bibinfo {author} {\bibfnamefont
  {S.}~\bibnamefont {Pagliara}},\ }\href {https://arxiv.org/abs/2412.02507}
  {\bibinfo {title} {Unveiling exciton formation: exploring the early stages in
  time, energy and momentum domain}} (\bibinfo {year} {2024}),\ \Eprint
  {https://arxiv.org/abs/2412.02507} {arXiv:2412.02507 [cond-mat.mtrl-sci]}
  \BibitemShut {NoStop}
\bibitem [{\citenamefont {Man}\ \emph {et~al.}(2021)\citenamefont {Man},
  \citenamefont {Madéo}, \citenamefont {Sahoo}, \citenamefont {Xie},
  \citenamefont {Campbell}, \citenamefont {Pareek}, \citenamefont {Karmakar},
  \citenamefont {Wong}, \citenamefont {Al-Mahboob}, \citenamefont {Chan},
  \citenamefont {Bacon}, \citenamefont {Zhu}, \citenamefont {Abdelrasoul},
  \citenamefont {Li}, \citenamefont {Heinz}, \citenamefont {da~Jornada},
  \citenamefont {Cao},\ and\ \citenamefont {Dani}}]{man_experimental_2021}
  \BibitemOpen
  \bibfield  {author} {\bibinfo {author} {\bibfnamefont {M.~K.~L.}\
  \bibnamefont {Man}}, \bibinfo {author} {\bibfnamefont {J.}~\bibnamefont
  {Madéo}}, \bibinfo {author} {\bibfnamefont {C.}~\bibnamefont {Sahoo}},
  \bibinfo {author} {\bibfnamefont {K.}~\bibnamefont {Xie}}, \bibinfo {author}
  {\bibfnamefont {M.}~\bibnamefont {Campbell}}, \bibinfo {author}
  {\bibfnamefont {V.}~\bibnamefont {Pareek}}, \bibinfo {author} {\bibfnamefont
  {A.}~\bibnamefont {Karmakar}}, \bibinfo {author} {\bibfnamefont {E.~L.}\
  \bibnamefont {Wong}}, \bibinfo {author} {\bibfnamefont {A.}~\bibnamefont
  {Al-Mahboob}}, \bibinfo {author} {\bibfnamefont {N.~S.}\ \bibnamefont
  {Chan}}, \bibinfo {author} {\bibfnamefont {D.~R.}\ \bibnamefont {Bacon}},
  \bibinfo {author} {\bibfnamefont {X.}~\bibnamefont {Zhu}}, \bibinfo {author}
  {\bibfnamefont {M.~M.~M.}\ \bibnamefont {Abdelrasoul}}, \bibinfo {author}
  {\bibfnamefont {X.}~\bibnamefont {Li}}, \bibinfo {author} {\bibfnamefont
  {T.~F.}\ \bibnamefont {Heinz}}, \bibinfo {author} {\bibfnamefont {F.~H.}\
  \bibnamefont {da~Jornada}}, \bibinfo {author} {\bibfnamefont
  {T.}~\bibnamefont {Cao}},\ and\ \bibinfo {author} {\bibfnamefont {K.~M.}\
  \bibnamefont {Dani}},\ }\bibfield  {title} {\bibinfo {title} {Experimental
  measurement of the intrinsic excitonic wave function},\ }\href
  {https://doi.org/10.1126/sciadv.abg0192} {\bibfield  {journal} {\bibinfo
  {journal} {Science Advances}\ }\textbf {\bibinfo {volume} {7}},\ \bibinfo
  {pages} {eabg0192} (\bibinfo {year} {2021})}\BibitemShut {NoStop}
\bibitem [{\citenamefont {Dong}\ \emph {et~al.}(2021)\citenamefont {Dong},
  \citenamefont {Puppin}, \citenamefont {Pincelli}, \citenamefont {Beaulieu},
  \citenamefont {Christiansen}, \citenamefont {Hübener}, \citenamefont
  {Nicholson}, \citenamefont {Xian}, \citenamefont {Dendzik}, \citenamefont
  {Deng}, \citenamefont {Windsor}, \citenamefont {Selig}, \citenamefont
  {Malic}, \citenamefont {Rubio}, \citenamefont {Knorr}, \citenamefont {Wolf},
  \citenamefont {Rettig},\ and\ \citenamefont {Ernstorfer}}]{dong_direct_2021}
  \BibitemOpen
  \bibfield  {author} {\bibinfo {author} {\bibfnamefont {S.}~\bibnamefont
  {Dong}}, \bibinfo {author} {\bibfnamefont {M.}~\bibnamefont {Puppin}},
  \bibinfo {author} {\bibfnamefont {T.}~\bibnamefont {Pincelli}}, \bibinfo
  {author} {\bibfnamefont {S.}~\bibnamefont {Beaulieu}}, \bibinfo {author}
  {\bibfnamefont {D.}~\bibnamefont {Christiansen}}, \bibinfo {author}
  {\bibfnamefont {H.}~\bibnamefont {Hübener}}, \bibinfo {author}
  {\bibfnamefont {C.~W.}\ \bibnamefont {Nicholson}}, \bibinfo {author}
  {\bibfnamefont {R.~P.}\ \bibnamefont {Xian}}, \bibinfo {author}
  {\bibfnamefont {M.}~\bibnamefont {Dendzik}}, \bibinfo {author} {\bibfnamefont
  {Y.}~\bibnamefont {Deng}}, \bibinfo {author} {\bibfnamefont {Y.~W.}\
  \bibnamefont {Windsor}}, \bibinfo {author} {\bibfnamefont {M.}~\bibnamefont
  {Selig}}, \bibinfo {author} {\bibfnamefont {E.}~\bibnamefont {Malic}},
  \bibinfo {author} {\bibfnamefont {A.}~\bibnamefont {Rubio}}, \bibinfo
  {author} {\bibfnamefont {A.}~\bibnamefont {Knorr}}, \bibinfo {author}
  {\bibfnamefont {M.}~\bibnamefont {Wolf}}, \bibinfo {author} {\bibfnamefont
  {L.}~\bibnamefont {Rettig}},\ and\ \bibinfo {author} {\bibfnamefont
  {R.}~\bibnamefont {Ernstorfer}},\ }\bibfield  {title} {\bibinfo {title}
  {Direct measurement of key exciton properties: Energy, dynamics, and spatial
  distribution of the wave function},\ }\href
  {https://doi.org/https://doi.org/10.1002/ntls.10010} {\bibfield  {journal}
  {\bibinfo  {journal} {Natural Sciences}\ }\textbf {\bibinfo {volume} {1}},\
  \bibinfo {pages} {e10010} (\bibinfo {year} {2021})}\BibitemShut {NoStop}
\bibitem [{\citenamefont {Meneghini}\ \emph {et~al.}(2023)\citenamefont
  {Meneghini}, \citenamefont {Reutzel}, \citenamefont {Mathias}, \citenamefont
  {Brem},\ and\ \citenamefont {Malic}}]{meneghini_hybrid_2023}
  \BibitemOpen
  \bibfield  {author} {\bibinfo {author} {\bibfnamefont {G.}~\bibnamefont
  {Meneghini}}, \bibinfo {author} {\bibfnamefont {M.}~\bibnamefont {Reutzel}},
  \bibinfo {author} {\bibfnamefont {S.}~\bibnamefont {Mathias}}, \bibinfo
  {author} {\bibfnamefont {S.}~\bibnamefont {Brem}},\ and\ \bibinfo {author}
  {\bibfnamefont {E.}~\bibnamefont {Malic}},\ }\bibfield  {title} {\bibinfo
  {title} {Hybrid exciton signatures in arpes spectra of van der waals
  materials},\ }\href {https://doi.org/10.1021/acsphotonics.3c00599} {\bibfield
   {journal} {\bibinfo  {journal} {ACS Photonics}\ }\textbf {\bibinfo {volume}
  {10}},\ \bibinfo {pages} {3570} (\bibinfo {year} {2023})}\BibitemShut
  {NoStop}
\bibitem [{\citenamefont {Mori}\ \emph {et~al.}(2023)\citenamefont {Mori},
  \citenamefont {Ciocys}, \citenamefont {Takasan}, \citenamefont {Ai},
  \citenamefont {Currier}, \citenamefont {Morimoto}, \citenamefont {Moore},\
  and\ \citenamefont {Lanzara}}]{mori_spin-polarized_2023}
  \BibitemOpen
  \bibfield  {author} {\bibinfo {author} {\bibfnamefont {R.}~\bibnamefont
  {Mori}}, \bibinfo {author} {\bibfnamefont {S.}~\bibnamefont {Ciocys}},
  \bibinfo {author} {\bibfnamefont {K.}~\bibnamefont {Takasan}}, \bibinfo
  {author} {\bibfnamefont {P.}~\bibnamefont {Ai}}, \bibinfo {author}
  {\bibfnamefont {K.}~\bibnamefont {Currier}}, \bibinfo {author} {\bibfnamefont
  {T.}~\bibnamefont {Morimoto}}, \bibinfo {author} {\bibfnamefont {J.~E.}\
  \bibnamefont {Moore}},\ and\ \bibinfo {author} {\bibfnamefont
  {A.}~\bibnamefont {Lanzara}},\ }\bibfield  {title} {\bibinfo {title}
  {Spin-polarized spatially indirect excitons in a topological insulator},\
  }\href {https://doi.org/10.1038/s41586-022-05567-3} {\bibfield  {journal}
  {\bibinfo  {journal} {Nature}\ }\textbf {\bibinfo {volume} {614}},\ \bibinfo
  {pages} {249} (\bibinfo {year} {2023})}\BibitemShut {NoStop}
\bibitem [{\citenamefont {Na}\ \emph {et~al.}(2023)\citenamefont {Na},
  \citenamefont {Mills},\ and\ \citenamefont {Jones}}]{na_advancing_2023}
  \BibitemOpen
  \bibfield  {author} {\bibinfo {author} {\bibfnamefont {M.}~\bibnamefont
  {Na}}, \bibinfo {author} {\bibfnamefont {A.~K.}\ \bibnamefont {Mills}},\ and\
  \bibinfo {author} {\bibfnamefont {D.~J.}\ \bibnamefont {Jones}},\ }\bibfield
  {title} {\bibinfo {title} {Advancing time- and angle-resolved photoemission
  spectroscopy: The role of ultrafast laser development},\ }\href
  {https://doi.org/https://doi.org/10.1016/j.physrep.2023.09.005} {\bibfield
  {journal} {\bibinfo  {journal} {Physics Reports}\ }\textbf {\bibinfo {volume}
  {1036}},\ \bibinfo {pages} {1} (\bibinfo {year} {2023})}\BibitemShut
  {NoStop}
\bibitem [{\citenamefont {Haug}\ and\ \citenamefont
  {Koch}(2009)}]{haug2009quantum}
  \BibitemOpen
  \bibfield  {author} {\bibinfo {author} {\bibfnamefont {H.}~\bibnamefont
  {Haug}}\ and\ \bibinfo {author} {\bibfnamefont {S.~W.}\ \bibnamefont
  {Koch}},\ }\href@noop {} {\emph {\bibinfo {title} {Quantum theory of the
  optical and electronic properties of semiconductors}}}\ (\bibinfo
  {publisher} {world scientific},\ \bibinfo {year} {2009})\BibitemShut
  {NoStop}
\bibitem [{\citenamefont {Lindberg}\ and\ \citenamefont
  {Koch}(1988)}]{lindberg_effective_1988}
  \BibitemOpen
  \bibfield  {author} {\bibinfo {author} {\bibfnamefont {M.}~\bibnamefont
  {Lindberg}}\ and\ \bibinfo {author} {\bibfnamefont {S.~W.}\ \bibnamefont
  {Koch}},\ }\bibfield  {title} {\bibinfo {title} {Effective bloch equations
  for semiconductors},\ }\href {https://doi.org/10.1103/PhysRevB.38.3342}
  {\bibfield  {journal} {\bibinfo  {journal} {Phys. Rev. B}\ }\textbf {\bibinfo
  {volume} {38}},\ \bibinfo {pages} {3342} (\bibinfo {year}
  {1988})}\BibitemShut {NoStop}
\bibitem [{\citenamefont {{\"O}streich}\ and\ \citenamefont
  {Sch{\"o}nhammer}(1993)}]{ostreich_non-stationary_1993}
  \BibitemOpen
  \bibfield  {author} {\bibinfo {author} {\bibfnamefont {T.}~\bibnamefont
  {{\"O}streich}}\ and\ \bibinfo {author} {\bibfnamefont {K.}~\bibnamefont
  {Sch{\"o}nhammer}},\ }\bibfield  {title} {\bibinfo {title} {Non-stationary
  excitonic-insulator states in photoexcited semiconductors},\ }\href
  {https://doi.org/10.1007/BF01315235} {\bibfield  {journal} {\bibinfo
  {journal} {Zeitschrift f{\"u}r Physik B Condensed Matter}\ }\textbf {\bibinfo
  {volume} {91}},\ \bibinfo {pages} {189} (\bibinfo {year} {1993})}\BibitemShut
  {NoStop}
\bibitem [{\citenamefont {Hannewald}\ \emph {et~al.}(2000)\citenamefont
  {Hannewald}, \citenamefont {Glutsch},\ and\ \citenamefont
  {Bechstedt}}]{hannewald_excitonic_2000}
  \BibitemOpen
  \bibfield  {author} {\bibinfo {author} {\bibfnamefont {K.}~\bibnamefont
  {Hannewald}}, \bibinfo {author} {\bibfnamefont {S.}~\bibnamefont {Glutsch}},\
  and\ \bibinfo {author} {\bibfnamefont {F.}~\bibnamefont {Bechstedt}},\
  }\bibfield  {title} {\bibinfo {title} {Excitonic insulator through coherent
  pulse excitation?},\ }\href
  {https://doi.org/10.1088\%2F0953-8984\%2F13\%2F2\%2F305} {\bibfield
  {journal} {\bibinfo  {journal} {Journal of Physics: Condensed Matter}\
  }\textbf {\bibinfo {volume} {13}},\ \bibinfo {pages} {275} (\bibinfo {year}
  {2000})}\BibitemShut {NoStop}
\bibitem [{\citenamefont {Koch}\ \emph {et~al.}(2006)\citenamefont {Koch},
  \citenamefont {Kira}, \citenamefont {Khitrova},\ and\ \citenamefont
  {Gibbs}}]{koch_semiconductor_2006}
  \BibitemOpen
  \bibfield  {author} {\bibinfo {author} {\bibfnamefont {S.~W.}\ \bibnamefont
  {Koch}}, \bibinfo {author} {\bibfnamefont {M.}~\bibnamefont {Kira}}, \bibinfo
  {author} {\bibfnamefont {G.}~\bibnamefont {Khitrova}},\ and\ \bibinfo
  {author} {\bibfnamefont {H.~M.}\ \bibnamefont {Gibbs}},\ }\bibfield  {title}
  {\bibinfo {title} {Semiconductor excitons in new light},\ }\href
  {https://doi.org/10.1038/nmat1658} {\bibfield  {journal} {\bibinfo  {journal}
  {Nature Materials}\ }\textbf {\bibinfo {volume} {5}},\ \bibinfo {pages} {523}
  (\bibinfo {year} {2006})}\BibitemShut {NoStop}
\bibitem [{\citenamefont {Madrid}\ \emph {et~al.}(2009)\citenamefont {Madrid},
  \citenamefont {Hyeon-Deuk}, \citenamefont {Habenicht},\ and\ \citenamefont
  {Prezhdo}}]{madrid_phonon-induced_2009}
  \BibitemOpen
  \bibfield  {author} {\bibinfo {author} {\bibfnamefont {A.~B.}\ \bibnamefont
  {Madrid}}, \bibinfo {author} {\bibfnamefont {K.}~\bibnamefont {Hyeon-Deuk}},
  \bibinfo {author} {\bibfnamefont {B.~F.}\ \bibnamefont {Habenicht}},\ and\
  \bibinfo {author} {\bibfnamefont {O.~V.}\ \bibnamefont {Prezhdo}},\
  }\bibfield  {title} {\bibinfo {title} {Phonon-induced dephasing of excitons
  in semiconductor quantum dots: Multiple exciton generation, fission, and
  luminescence},\ }\href {https://doi.org/10.1021/nn900584p} {\bibfield
  {journal} {\bibinfo  {journal} {ACS Nano}\ }\textbf {\bibinfo {volume} {3}},\
  \bibinfo {pages} {2487} (\bibinfo {year} {2009})}\BibitemShut {NoStop}
\bibitem [{\citenamefont {Nie}\ \emph {et~al.}(2014)\citenamefont {Nie},
  \citenamefont {Long}, \citenamefont {Sun}, \citenamefont {Huang},
  \citenamefont {Zhang}, \citenamefont {Xiong}, \citenamefont {Hewak},
  \citenamefont {Shen}, \citenamefont {Prezhdo},\ and\ \citenamefont
  {Loh}}]{nie_ultrafast_2014}
  \BibitemOpen
  \bibfield  {author} {\bibinfo {author} {\bibfnamefont {Z.}~\bibnamefont
  {Nie}}, \bibinfo {author} {\bibfnamefont {R.}~\bibnamefont {Long}}, \bibinfo
  {author} {\bibfnamefont {L.}~\bibnamefont {Sun}}, \bibinfo {author}
  {\bibfnamefont {C.-C.}\ \bibnamefont {Huang}}, \bibinfo {author}
  {\bibfnamefont {J.}~\bibnamefont {Zhang}}, \bibinfo {author} {\bibfnamefont
  {Q.}~\bibnamefont {Xiong}}, \bibinfo {author} {\bibfnamefont {D.~W.}\
  \bibnamefont {Hewak}}, \bibinfo {author} {\bibfnamefont {Z.}~\bibnamefont
  {Shen}}, \bibinfo {author} {\bibfnamefont {O.~V.}\ \bibnamefont {Prezhdo}},\
  and\ \bibinfo {author} {\bibfnamefont {Z.-H.}\ \bibnamefont {Loh}},\
  }\bibfield  {title} {\bibinfo {title} {Ultrafast carrier thermalization and
  cooling dynamics in few-layer {M}o{S}$_2$},\ }\href
  {https://doi.org/10.1021/nn504760x} {\bibfield  {journal} {\bibinfo
  {journal} {ACS Nano}\ }\textbf {\bibinfo {volume} {8}},\ \bibinfo {pages}
  {10931} (\bibinfo {year} {2014})}\BibitemShut {NoStop}
\bibitem [{\citenamefont {Dey}\ \emph {et~al.}(2016)\citenamefont {Dey},
  \citenamefont {Paul}, \citenamefont {Wang}, \citenamefont {Stevens},
  \citenamefont {Liu}, \citenamefont {Romero}, \citenamefont {Shan},
  \citenamefont {Hilton},\ and\ \citenamefont {Karaiskaj}}]{dey_optical_2016}
  \BibitemOpen
  \bibfield  {author} {\bibinfo {author} {\bibfnamefont {P.}~\bibnamefont
  {Dey}}, \bibinfo {author} {\bibfnamefont {J.}~\bibnamefont {Paul}}, \bibinfo
  {author} {\bibfnamefont {Z.}~\bibnamefont {Wang}}, \bibinfo {author}
  {\bibfnamefont {C.~E.}\ \bibnamefont {Stevens}}, \bibinfo {author}
  {\bibfnamefont {C.}~\bibnamefont {Liu}}, \bibinfo {author} {\bibfnamefont
  {A.~H.}\ \bibnamefont {Romero}}, \bibinfo {author} {\bibfnamefont
  {J.}~\bibnamefont {Shan}}, \bibinfo {author} {\bibfnamefont {D.~J.}\
  \bibnamefont {Hilton}},\ and\ \bibinfo {author} {\bibfnamefont
  {D.}~\bibnamefont {Karaiskaj}},\ }\bibfield  {title} {\bibinfo {title}
  {Optical coherence in atomic-monolayer transition-metal dichalcogenides
  limited by electron-phonon interactions},\ }\href
  {https://doi.org/10.1103/PhysRevLett.116.127402} {\bibfield  {journal}
  {\bibinfo  {journal} {Phys. Rev. Lett.}\ }\textbf {\bibinfo {volume} {116}},\
  \bibinfo {pages} {127402} (\bibinfo {year} {2016})}\BibitemShut {NoStop}
\bibitem [{\citenamefont {Sangalli}\ \emph {et~al.}(2018)\citenamefont
  {Sangalli}, \citenamefont {Perfetto}, \citenamefont {Stefanucci},\ and\
  \citenamefont {Marini}}]{sangalli_an-ab-initio_2018}
  \BibitemOpen
  \bibfield  {author} {\bibinfo {author} {\bibfnamefont {D.}~\bibnamefont
  {Sangalli}}, \bibinfo {author} {\bibfnamefont {E.}~\bibnamefont {Perfetto}},
  \bibinfo {author} {\bibfnamefont {G.}~\bibnamefont {Stefanucci}},\ and\
  \bibinfo {author} {\bibfnamefont {A.}~\bibnamefont {Marini}},\ }\bibfield
  {title} {\bibinfo {title} {An ab-initio approach to describe coherent and
  non-coherent exciton dynamics},\ }\href
  {https://doi.org/10.1140/epjb/e2018-90126-5} {\bibfield  {journal} {\bibinfo
  {journal} {The European Physical Journal B}\ }\textbf {\bibinfo {volume}
  {91}},\ \bibinfo {pages} {171} (\bibinfo {year} {2018})}\BibitemShut
  {NoStop}
\bibitem [{\citenamefont {Stefanucci}\ and\ \citenamefont
  {Perfetto}(2021)}]{stefanucci_from-carriers_2021}
  \BibitemOpen
  \bibfield  {author} {\bibinfo {author} {\bibfnamefont {G.}~\bibnamefont
  {Stefanucci}}\ and\ \bibinfo {author} {\bibfnamefont {E.}~\bibnamefont
  {Perfetto}},\ }\bibfield  {title} {\bibinfo {title} {From carriers and
  virtual excitons to exciton populations: Insights into time-resolved arpes
  spectra from an exactly solvable model},\ }\href
  {https://doi.org/10.1103/PhysRevB.103.245103} {\bibfield  {journal} {\bibinfo
   {journal} {Phys. Rev. B}\ }\textbf {\bibinfo {volume} {103}},\ \bibinfo
  {pages} {245103} (\bibinfo {year} {2021})}\BibitemShut {NoStop}
\bibitem [{\citenamefont {Perfetto}\ and\ \citenamefont
  {Stefanucci}(2023)}]{perfetto_real-time_2023}
  \BibitemOpen
  \bibfield  {author} {\bibinfo {author} {\bibfnamefont {E.}~\bibnamefont
  {Perfetto}}\ and\ \bibinfo {author} {\bibfnamefont {G.}~\bibnamefont
  {Stefanucci}},\ }\bibfield  {title} {\bibinfo {title} {Real-time
  {G}{W}-{E}hrenfest-{F}an-{M}igdal method for nonequilibrium 2d materials},\
  }\href {https://doi.org/10.1021/acs.nanolett.3c01772} {\bibfield  {journal}
  {\bibinfo  {journal} {Nano Letters}\ }\textbf {\bibinfo {volume} {23}},\
  \bibinfo {pages} {7029} (\bibinfo {year} {2023})}\BibitemShut {NoStop}
\bibitem [{\citenamefont {Bertoni}\ \emph {et~al.}(2016)\citenamefont
  {Bertoni}, \citenamefont {Nicholson}, \citenamefont {Waldecker},
  \citenamefont {H\"ubener}, \citenamefont {Monney}, \citenamefont
  {De~Giovannini}, \citenamefont {Puppin}, \citenamefont {Hoesch},
  \citenamefont {Springate}, \citenamefont {Chapman}, \citenamefont {Cacho},
  \citenamefont {Wolf}, \citenamefont {Rubio},\ and\ \citenamefont
  {Ernstorfer}}]{bertoni_generation_2016}
  \BibitemOpen
  \bibfield  {author} {\bibinfo {author} {\bibfnamefont {R.}~\bibnamefont
  {Bertoni}}, \bibinfo {author} {\bibfnamefont {C.~W.}\ \bibnamefont
  {Nicholson}}, \bibinfo {author} {\bibfnamefont {L.}~\bibnamefont
  {Waldecker}}, \bibinfo {author} {\bibfnamefont {H.}~\bibnamefont
  {H\"ubener}}, \bibinfo {author} {\bibfnamefont {C.}~\bibnamefont {Monney}},
  \bibinfo {author} {\bibfnamefont {U.}~\bibnamefont {De~Giovannini}}, \bibinfo
  {author} {\bibfnamefont {M.}~\bibnamefont {Puppin}}, \bibinfo {author}
  {\bibfnamefont {M.}~\bibnamefont {Hoesch}}, \bibinfo {author} {\bibfnamefont
  {E.}~\bibnamefont {Springate}}, \bibinfo {author} {\bibfnamefont {R.~T.}\
  \bibnamefont {Chapman}}, \bibinfo {author} {\bibfnamefont {C.}~\bibnamefont
  {Cacho}}, \bibinfo {author} {\bibfnamefont {M.}~\bibnamefont {Wolf}},
  \bibinfo {author} {\bibfnamefont {A.}~\bibnamefont {Rubio}},\ and\ \bibinfo
  {author} {\bibfnamefont {R.}~\bibnamefont {Ernstorfer}},\ }\bibfield  {title}
  {\bibinfo {title} {Generation and evolution of spin-, valley-, and
  layer-polarized excited carriers in inversion-symmetric {WS}e$_{2}$},\ }\href
  {https://doi.org/10.1103/PhysRevLett.117.277201} {\bibfield  {journal}
  {\bibinfo  {journal} {Phys. Rev. Lett.}\ }\textbf {\bibinfo {volume} {117}},\
  \bibinfo {pages} {277201} (\bibinfo {year} {2016})}\BibitemShut {NoStop}
\bibitem [{\citenamefont {Uddin}\ \emph {et~al.}(2020)\citenamefont {Uddin},
  \citenamefont {Kim}, \citenamefont {Lorenzon}, \citenamefont {Yeh},
  \citenamefont {Lien}, \citenamefont {Barnard}, \citenamefont {Htoon},
  \citenamefont {Weber-Bargioni},\ and\ \citenamefont
  {Javey}}]{uddin_neutral_2020}
  \BibitemOpen
  \bibfield  {author} {\bibinfo {author} {\bibfnamefont {S.~Z.}\ \bibnamefont
  {Uddin}}, \bibinfo {author} {\bibfnamefont {H.}~\bibnamefont {Kim}}, \bibinfo
  {author} {\bibfnamefont {M.}~\bibnamefont {Lorenzon}}, \bibinfo {author}
  {\bibfnamefont {M.}~\bibnamefont {Yeh}}, \bibinfo {author} {\bibfnamefont
  {D.-H.}\ \bibnamefont {Lien}}, \bibinfo {author} {\bibfnamefont {E.~S.}\
  \bibnamefont {Barnard}}, \bibinfo {author} {\bibfnamefont {H.}~\bibnamefont
  {Htoon}}, \bibinfo {author} {\bibfnamefont {A.}~\bibnamefont
  {Weber-Bargioni}},\ and\ \bibinfo {author} {\bibfnamefont {A.}~\bibnamefont
  {Javey}},\ }\bibfield  {title} {\bibinfo {title} {Neutral exciton diffusion
  in monolayer {M}o{S}$_2$},\ }\href {https://doi.org/10.1021/acsnano.0c05305}
  {\bibfield  {journal} {\bibinfo  {journal} {ACS Nano}\ }\textbf {\bibinfo
  {volume} {14}},\ \bibinfo {pages} {13433} (\bibinfo {year}
  {2020})}\BibitemShut {NoStop}
\bibitem [{\citenamefont {Perfetto}\ and\ \citenamefont
  {Stefanucci}(2021)}]{perfetto_ultrafast_2021}
  \BibitemOpen
  \bibfield  {author} {\bibinfo {author} {\bibfnamefont {E.}~\bibnamefont
  {Perfetto}}\ and\ \bibinfo {author} {\bibfnamefont {G.}~\bibnamefont
  {Stefanucci}},\ }\bibfield  {title} {\bibinfo {title} {Ultrafast creation and
  melting of nonequilibrium excitonic condensates in bulk {WS}e$_{2}$},\ }\href
  {https://doi.org/10.1103/PhysRevB.103.L241404} {\bibfield  {journal}
  {\bibinfo  {journal} {Phys. Rev. B}\ }\textbf {\bibinfo {volume} {103}},\
  \bibinfo {pages} {L241404} (\bibinfo {year} {2021})}\BibitemShut {NoStop}
\bibitem [{\citenamefont {Stefanucci}\ and\ \citenamefont
  {Perfetto}(2025)}]{stefanucci_excitonic_2025}
  \BibitemOpen
  \bibfield  {author} {\bibinfo {author} {\bibfnamefont {G.}~\bibnamefont
  {Stefanucci}}\ and\ \bibinfo {author} {\bibfnamefont {E.}~\bibnamefont
  {Perfetto}},\ }\bibfield  {title} {\bibinfo {title} {{Excitonic Bloch
  equations from first principles}},\ }\href
  {https://doi.org/10.21468/SciPostPhys.18.1.009} {\bibfield  {journal}
  {\bibinfo  {journal} {SciPost Phys.}\ }\textbf {\bibinfo {volume} {18}},\
  \bibinfo {pages} {009} (\bibinfo {year} {2025})}\BibitemShut {NoStop}
\bibitem [{\citenamefont {Moody}\ \emph {et~al.}(2015)\citenamefont {Moody},
  \citenamefont {Kavir~Dass}, \citenamefont {Hao}, \citenamefont {Chen},
  \citenamefont {Li}, \citenamefont {Singh}, \citenamefont {Tran},
  \citenamefont {Clark}, \citenamefont {Xu}, \citenamefont {Bergh{\"a}user},
  \citenamefont {Malic}, \citenamefont {Knorr},\ and\ \citenamefont
  {Li}}]{moody_intrinsic_2015}
  \BibitemOpen
  \bibfield  {author} {\bibinfo {author} {\bibfnamefont {G.}~\bibnamefont
  {Moody}}, \bibinfo {author} {\bibfnamefont {C.}~\bibnamefont {Kavir~Dass}},
  \bibinfo {author} {\bibfnamefont {K.}~\bibnamefont {Hao}}, \bibinfo {author}
  {\bibfnamefont {C.-H.}\ \bibnamefont {Chen}}, \bibinfo {author}
  {\bibfnamefont {L.-J.}\ \bibnamefont {Li}}, \bibinfo {author} {\bibfnamefont
  {A.}~\bibnamefont {Singh}}, \bibinfo {author} {\bibfnamefont
  {K.}~\bibnamefont {Tran}}, \bibinfo {author} {\bibfnamefont {G.}~\bibnamefont
  {Clark}}, \bibinfo {author} {\bibfnamefont {X.}~\bibnamefont {Xu}}, \bibinfo
  {author} {\bibfnamefont {G.}~\bibnamefont {Bergh{\"a}user}}, \bibinfo
  {author} {\bibfnamefont {E.}~\bibnamefont {Malic}}, \bibinfo {author}
  {\bibfnamefont {A.}~\bibnamefont {Knorr}},\ and\ \bibinfo {author}
  {\bibfnamefont {X.}~\bibnamefont {Li}},\ }\bibfield  {title} {\bibinfo
  {title} {Intrinsic homogeneous linewidth and broadening mechanisms of
  excitons in monolayer transition metal dichalcogenides},\ }\href
  {https://doi.org/10.1038/ncomms9315} {\bibfield  {journal} {\bibinfo
  {journal} {Nature Communications}\ }\textbf {\bibinfo {volume} {6}},\
  \bibinfo {pages} {8315} (\bibinfo {year} {2015})}\BibitemShut {NoStop}
\bibitem [{\citenamefont {Selig}\ \emph {et~al.}(2016)\citenamefont {Selig},
  \citenamefont {Bergh{\"a}user}, \citenamefont {Raja}, \citenamefont {Nagler},
  \citenamefont {Sch{\"u}ller}, \citenamefont {Heinz}, \citenamefont {Korn},
  \citenamefont {Chernikov}, \citenamefont {Malic},\ and\ \citenamefont
  {Knorr}}]{selig_exciton_2016}
  \BibitemOpen
  \bibfield  {author} {\bibinfo {author} {\bibfnamefont {M.}~\bibnamefont
  {Selig}}, \bibinfo {author} {\bibfnamefont {G.}~\bibnamefont
  {Bergh{\"a}user}}, \bibinfo {author} {\bibfnamefont {A.}~\bibnamefont
  {Raja}}, \bibinfo {author} {\bibfnamefont {P.}~\bibnamefont {Nagler}},
  \bibinfo {author} {\bibfnamefont {C.}~\bibnamefont {Sch{\"u}ller}}, \bibinfo
  {author} {\bibfnamefont {T.~F.}\ \bibnamefont {Heinz}}, \bibinfo {author}
  {\bibfnamefont {T.}~\bibnamefont {Korn}}, \bibinfo {author} {\bibfnamefont
  {A.}~\bibnamefont {Chernikov}}, \bibinfo {author} {\bibfnamefont
  {E.}~\bibnamefont {Malic}},\ and\ \bibinfo {author} {\bibfnamefont
  {A.}~\bibnamefont {Knorr}},\ }\bibfield  {title} {\bibinfo {title} {Excitonic
  linewidth and coherence lifetime in monolayer transition metal
  dichalcogenides},\ }\href {https://doi.org/10.1038/ncomms13279} {\bibfield
  {journal} {\bibinfo  {journal} {Nature Communications}\ }\textbf {\bibinfo
  {volume} {7}},\ \bibinfo {pages} {13279} (\bibinfo {year}
  {2016})}\BibitemShut {NoStop}
\bibitem [{\citenamefont {Jakubczyk}\ \emph {et~al.}(2018)\citenamefont
  {Jakubczyk}, \citenamefont {Nogajewski}, \citenamefont {Molas}, \citenamefont
  {Bartos}, \citenamefont {Langbein}, \citenamefont {Potemski},\ and\
  \citenamefont {Kasprzak}}]{jakubczyk_impact_2018}
  \BibitemOpen
  \bibfield  {author} {\bibinfo {author} {\bibfnamefont {T.}~\bibnamefont
  {Jakubczyk}}, \bibinfo {author} {\bibfnamefont {K.}~\bibnamefont
  {Nogajewski}}, \bibinfo {author} {\bibfnamefont {M.~R.}\ \bibnamefont
  {Molas}}, \bibinfo {author} {\bibfnamefont {M.}~\bibnamefont {Bartos}},
  \bibinfo {author} {\bibfnamefont {W.}~\bibnamefont {Langbein}}, \bibinfo
  {author} {\bibfnamefont {M.}~\bibnamefont {Potemski}},\ and\ \bibinfo
  {author} {\bibfnamefont {J.}~\bibnamefont {Kasprzak}},\ }\bibfield  {title}
  {\bibinfo {title} {Impact of environment on dynamics of exciton complexes in
  a ws2 monolayer},\ }\href {https://doi.org/10.1088/2053-1583/aabc1c}
  {\bibfield  {journal} {\bibinfo  {journal} {2D Materials}\ }\textbf {\bibinfo
  {volume} {5}},\ \bibinfo {pages} {031007} (\bibinfo {year}
  {2018})}\BibitemShut {NoStop}
\bibitem [{\citenamefont {Perfetto}\ \emph {et~al.}(2019)\citenamefont
  {Perfetto}, \citenamefont {Sangalli}, \citenamefont {Marini},\ and\
  \citenamefont {Stefanucci}}]{perfetto_pump-driven_2019}
  \BibitemOpen
  \bibfield  {author} {\bibinfo {author} {\bibfnamefont {E.}~\bibnamefont
  {Perfetto}}, \bibinfo {author} {\bibfnamefont {D.}~\bibnamefont {Sangalli}},
  \bibinfo {author} {\bibfnamefont {A.}~\bibnamefont {Marini}},\ and\ \bibinfo
  {author} {\bibfnamefont {G.}~\bibnamefont {Stefanucci}},\ }\bibfield  {title}
  {\bibinfo {title} {Pump-driven normal-to-excitonic insulator transition:
  Josephson oscillations and signatures of {BEC-BCS} crossover in time-resolved
  {ARPES}},\ }\href {https://doi.org/10.1103/PhysRevMaterials.3.124601}
  {\bibfield  {journal} {\bibinfo  {journal} {Phys. Rev. Materials}\ }\textbf
  {\bibinfo {volume} {3}},\ \bibinfo {pages} {124601} (\bibinfo {year}
  {2019})}\BibitemShut {NoStop}
\bibitem [{\citenamefont {Perfetto}\ \emph {et~al.}(2020)\citenamefont
  {Perfetto}, \citenamefont {Bianchi},\ and\ \citenamefont
  {Stefanucci}}]{perfetto_time-resolved_2020}
  \BibitemOpen
  \bibfield  {author} {\bibinfo {author} {\bibfnamefont {E.}~\bibnamefont
  {Perfetto}}, \bibinfo {author} {\bibfnamefont {S.}~\bibnamefont {Bianchi}},\
  and\ \bibinfo {author} {\bibfnamefont {G.}~\bibnamefont {Stefanucci}},\
  }\bibfield  {title} {\bibinfo {title} {Time-resolved arpes spectra of
  nonequilibrium excitonic insulators: Revealing macroscopic coherence with
  ultrashort pulses},\ }\href {https://doi.org/10.1103/PhysRevB.101.041201}
  {\bibfield  {journal} {\bibinfo  {journal} {Phys. Rev. B}\ }\textbf {\bibinfo
  {volume} {101}},\ \bibinfo {pages} {041201} (\bibinfo {year}
  {2020})}\BibitemShut {NoStop}
\bibitem [{\citenamefont {Perfetto}\ and\ \citenamefont
  {Stefanucci}(2020)}]{perfetto_floquet_2020}
  \BibitemOpen
  \bibfield  {author} {\bibinfo {author} {\bibfnamefont {E.}~\bibnamefont
  {Perfetto}}\ and\ \bibinfo {author} {\bibfnamefont {G.}~\bibnamefont
  {Stefanucci}},\ }\bibfield  {title} {\bibinfo {title} {Floquet topological
  phase of nondriven $p$-wave nonequilibrium excitonic insulators},\ }\href
  {https://link.aps.org/doi/10.1103/PhysRevLett.125.106401} {\bibfield
  {journal} {\bibinfo  {journal} {Phys. Rev. Lett.}\ }\textbf {\bibinfo
  {volume} {125}},\ \bibinfo {pages} {106401} (\bibinfo {year}
  {2020})}\BibitemShut {NoStop}
\bibitem [{\citenamefont {Chan}\ \emph {et~al.}(2023)\citenamefont {Chan},
  \citenamefont {Qiu}, \citenamefont {da~Jornada},\ and\ \citenamefont
  {Louie}}]{chan_giant_2023}
  \BibitemOpen
  \bibfield  {author} {\bibinfo {author} {\bibfnamefont {Y.-H.}\ \bibnamefont
  {Chan}}, \bibinfo {author} {\bibfnamefont {D.~Y.}\ \bibnamefont {Qiu}},
  \bibinfo {author} {\bibfnamefont {F.~H.}\ \bibnamefont {da~Jornada}},\ and\
  \bibinfo {author} {\bibfnamefont {S.~G.}\ \bibnamefont {Louie}},\ }\bibfield
  {title} {\bibinfo {title} {Giant self-driven exciton-floquet signatures in
  time-resolved photoemission spectroscopy of {MoS$_2$} from time-dependent
  {GW} approach},\ }\href {https://doi.org/10.1073/pnas.230195712} {\bibfield
  {journal} {\bibinfo  {journal} {Proceedings of the National Academy of
  Sciences}\ }\textbf {\bibinfo {volume} {120}},\ \bibinfo {pages}
  {e2301957120} (\bibinfo {year} {2023})}\BibitemShut {NoStop}
\bibitem [{\citenamefont {Pareek}\ \emph {et~al.}(2026)\citenamefont {Pareek},
  \citenamefont {Bacon}, \citenamefont {Zhu}, \citenamefont {Chan},
  \citenamefont {Bussolotti}, \citenamefont {Menezes}, \citenamefont {Chan},
  \citenamefont {Urquizo}, \citenamefont {Watanabe}, \citenamefont {Taniguchi},
  \citenamefont {Perfetto}, \citenamefont {Man}, \citenamefont {Mad{\'e}o},
  \citenamefont {Stefanucci}, \citenamefont {Qiu}, \citenamefont {Goh},
  \citenamefont {da~Jornada},\ and\ \citenamefont
  {Dani}}]{pareek_driving_2026}
  \BibitemOpen
  \bibfield  {author} {\bibinfo {author} {\bibfnamefont {V.}~\bibnamefont
  {Pareek}}, \bibinfo {author} {\bibfnamefont {D.~R.}\ \bibnamefont {Bacon}},
  \bibinfo {author} {\bibfnamefont {X.}~\bibnamefont {Zhu}}, \bibinfo {author}
  {\bibfnamefont {Y.-H.}\ \bibnamefont {Chan}}, \bibinfo {author}
  {\bibfnamefont {F.}~\bibnamefont {Bussolotti}}, \bibinfo {author}
  {\bibfnamefont {M.~G.}\ \bibnamefont {Menezes}}, \bibinfo {author}
  {\bibfnamefont {N.~S.}\ \bibnamefont {Chan}}, \bibinfo {author}
  {\bibfnamefont {J.~P.}\ \bibnamefont {Urquizo}}, \bibinfo {author}
  {\bibfnamefont {K.}~\bibnamefont {Watanabe}}, \bibinfo {author}
  {\bibfnamefont {T.}~\bibnamefont {Taniguchi}}, \bibinfo {author}
  {\bibfnamefont {E.}~\bibnamefont {Perfetto}}, \bibinfo {author}
  {\bibfnamefont {M.~K.~L.}\ \bibnamefont {Man}}, \bibinfo {author}
  {\bibfnamefont {J.}~\bibnamefont {Mad{\'e}o}}, \bibinfo {author}
  {\bibfnamefont {G.}~\bibnamefont {Stefanucci}}, \bibinfo {author}
  {\bibfnamefont {D.~Y.}\ \bibnamefont {Qiu}}, \bibinfo {author} {\bibfnamefont
  {K.~E.~J.}\ \bibnamefont {Goh}}, \bibinfo {author} {\bibfnamefont {F.~H.}\
  \bibnamefont {da~Jornada}},\ and\ \bibinfo {author} {\bibfnamefont {K.~M.}\
  \bibnamefont {Dani}},\ }\bibfield  {title} {\bibinfo {title} {Driving floquet
  physics with excitonic fields},\ }\bibfield  {journal} {\bibinfo  {journal}
  {Nature Physics}\ }\href {https://doi.org/10.1038/s41567-025-03132-z}
  {10.1038/s41567-025-03132-z} (\bibinfo {year} {2026})\BibitemShut {NoStop}
\bibitem [{\citenamefont {Perfetto}\ \emph {et~al.}(2016)\citenamefont
  {Perfetto}, \citenamefont {Sangalli}, \citenamefont {Marini},\ and\
  \citenamefont {Stefanucci}}]{perfetto_first-principles_2016}
  \BibitemOpen
  \bibfield  {author} {\bibinfo {author} {\bibfnamefont {E.}~\bibnamefont
  {Perfetto}}, \bibinfo {author} {\bibfnamefont {D.}~\bibnamefont {Sangalli}},
  \bibinfo {author} {\bibfnamefont {A.}~\bibnamefont {Marini}},\ and\ \bibinfo
  {author} {\bibfnamefont {G.}~\bibnamefont {Stefanucci}},\ }\bibfield  {title}
  {\bibinfo {title} {First-principles approach to excitons in time-resolved and
  angle-resolved photoemission spectra},\ }\href
  {https://doi.org/10.1103/PhysRevB.94.245303} {\bibfield  {journal} {\bibinfo
  {journal} {Phys. Rev. B}\ }\textbf {\bibinfo {volume} {94}},\ \bibinfo
  {pages} {245303} (\bibinfo {year} {2016})}\BibitemShut {NoStop}
\bibitem [{\citenamefont {Steinhoff}\ \emph {et~al.}(2017)\citenamefont
  {Steinhoff}, \citenamefont {Florian}, \citenamefont {R{\"o}sner},
  \citenamefont {Sch{\"o}nhoff}, \citenamefont {Wehling},\ and\ \citenamefont
  {Jahnke}}]{steinhoff_exciton_2017}
  \BibitemOpen
  \bibfield  {author} {\bibinfo {author} {\bibfnamefont {A.}~\bibnamefont
  {Steinhoff}}, \bibinfo {author} {\bibfnamefont {M.}~\bibnamefont {Florian}},
  \bibinfo {author} {\bibfnamefont {M.}~\bibnamefont {R{\"o}sner}}, \bibinfo
  {author} {\bibfnamefont {G.}~\bibnamefont {Sch{\"o}nhoff}}, \bibinfo {author}
  {\bibfnamefont {T.~O.}\ \bibnamefont {Wehling}},\ and\ \bibinfo {author}
  {\bibfnamefont {F.}~\bibnamefont {Jahnke}},\ }\bibfield  {title} {\bibinfo
  {title} {Exciton fission in monolayer transition metal dichalcogenide
  semiconductors},\ }\href {https://doi.org/10.1038/s41467-017-01298-6}
  {\bibfield  {journal} {\bibinfo  {journal} {Nature Communications}\ }\textbf
  {\bibinfo {volume} {8}},\ \bibinfo {pages} {1166} (\bibinfo {year}
  {2017})}\BibitemShut {NoStop}
\bibitem [{\citenamefont {Rustagi}\ and\ \citenamefont
  {Kemper}(2018)}]{rustagi_photoemission_2018}
  \BibitemOpen
  \bibfield  {author} {\bibinfo {author} {\bibfnamefont {A.}~\bibnamefont
  {Rustagi}}\ and\ \bibinfo {author} {\bibfnamefont {A.~F.}\ \bibnamefont
  {Kemper}},\ }\bibfield  {title} {\bibinfo {title} {Photoemission signature of
  excitons},\ }\href {https://doi.org/10.1103/PhysRevB.97.235310} {\bibfield
  {journal} {\bibinfo  {journal} {Phys. Rev. B}\ }\textbf {\bibinfo {volume}
  {97}},\ \bibinfo {pages} {235310} (\bibinfo {year} {2018})}\BibitemShut
  {NoStop}
\bibitem [{\citenamefont {Christiansen}\ \emph {et~al.}(2019)\citenamefont
  {Christiansen}, \citenamefont {Selig}, \citenamefont {Malic}, \citenamefont
  {Ernstorfer},\ and\ \citenamefont {Knorr}}]{christiansen_theory_2019}
  \BibitemOpen
  \bibfield  {author} {\bibinfo {author} {\bibfnamefont {D.}~\bibnamefont
  {Christiansen}}, \bibinfo {author} {\bibfnamefont {M.}~\bibnamefont {Selig}},
  \bibinfo {author} {\bibfnamefont {E.}~\bibnamefont {Malic}}, \bibinfo
  {author} {\bibfnamefont {R.}~\bibnamefont {Ernstorfer}},\ and\ \bibinfo
  {author} {\bibfnamefont {A.}~\bibnamefont {Knorr}},\ }\bibfield  {title}
  {\bibinfo {title} {Theory of exciton dynamics in time-resolved arpes: Intra-
  and intervalley scattering in two-dimensional semiconductors},\ }\href
  {https://doi.org/10.1103/PhysRevB.100.205401} {\bibfield  {journal} {\bibinfo
   {journal} {Phys. Rev. B}\ }\textbf {\bibinfo {volume} {100}},\ \bibinfo
  {pages} {205401} (\bibinfo {year} {2019})}\BibitemShut {NoStop}
\bibitem [{\citenamefont {Freericks}\ \emph {et~al.}(2009)\citenamefont
  {Freericks}, \citenamefont {Krishnamurthy},\ and\ \citenamefont
  {Pruschke}}]{freericks_theoretical_2009}
  \BibitemOpen
  \bibfield  {author} {\bibinfo {author} {\bibfnamefont {J.~K.}\ \bibnamefont
  {Freericks}}, \bibinfo {author} {\bibfnamefont {H.~R.}\ \bibnamefont
  {Krishnamurthy}},\ and\ \bibinfo {author} {\bibfnamefont {T.}~\bibnamefont
  {Pruschke}},\ }\bibfield  {title} {\bibinfo {title} {Theoretical description
  of time-resolved photoemission spectroscopy: Application to pump-probe
  experiments},\ }\href {https://doi.org/10.1103/PhysRevLett.102.136401}
  {\bibfield  {journal} {\bibinfo  {journal} {Phys. Rev. Lett.}\ }\textbf
  {\bibinfo {volume} {102}},\ \bibinfo {pages} {136401} (\bibinfo {year}
  {2009})}\BibitemShut {NoStop}
\bibitem [{\citenamefont {Stefanucci}\ and\ \citenamefont {van
  Leeuwen}(2025)}]{svl-book_2025}
  \BibitemOpen
  \bibfield  {author} {\bibinfo {author} {\bibfnamefont {G.}~\bibnamefont
  {Stefanucci}}\ and\ \bibinfo {author} {\bibfnamefont {R.}~\bibnamefont {van
  Leeuwen}},\ }\href@noop {} {\emph {\bibinfo {title} {Nonequilibrium Many-Body
  Theory of Quantum Systems: A Modern Introduction}}},\ \bibinfo {edition}
  {2nd}\ ed.\ (\bibinfo  {publisher} {Cambridge University Press},\ \bibinfo
  {year} {2025})\BibitemShut {NoStop}
\bibitem [{\citenamefont {Thr\"anhardt}\ \emph {et~al.}(2000)\citenamefont
  {Thr\"anhardt}, \citenamefont {Kuckenburg}, \citenamefont {Knorr},
  \citenamefont {Meier},\ and\ \citenamefont {Koch}}]{thranhardt_quantum_2000}
  \BibitemOpen
  \bibfield  {author} {\bibinfo {author} {\bibfnamefont {A.}~\bibnamefont
  {Thr\"anhardt}}, \bibinfo {author} {\bibfnamefont {S.}~\bibnamefont
  {Kuckenburg}}, \bibinfo {author} {\bibfnamefont {A.}~\bibnamefont {Knorr}},
  \bibinfo {author} {\bibfnamefont {T.}~\bibnamefont {Meier}},\ and\ \bibinfo
  {author} {\bibfnamefont {S.~W.}\ \bibnamefont {Koch}},\ }\bibfield  {title}
  {\bibinfo {title} {Quantum theory of phonon-assisted exciton formation and
  luminescence in semiconductor quantum wells},\ }\href
  {https://doi.org/10.1103/PhysRevB.62.2706} {\bibfield  {journal} {\bibinfo
  {journal} {Phys. Rev. B}\ }\textbf {\bibinfo {volume} {62}},\ \bibinfo
  {pages} {2706} (\bibinfo {year} {2000})}\BibitemShut {NoStop}
\bibitem [{\citenamefont {Brem}\ \emph {et~al.}(2018)\citenamefont {Brem},
  \citenamefont {Selig}, \citenamefont {Berghaeuser},\ and\ \citenamefont
  {Malic}}]{brem_exciton_2018}
  \BibitemOpen
  \bibfield  {author} {\bibinfo {author} {\bibfnamefont {S.}~\bibnamefont
  {Brem}}, \bibinfo {author} {\bibfnamefont {M.}~\bibnamefont {Selig}},
  \bibinfo {author} {\bibfnamefont {G.}~\bibnamefont {Berghaeuser}},\ and\
  \bibinfo {author} {\bibfnamefont {E.}~\bibnamefont {Malic}},\ }\bibfield
  {title} {\bibinfo {title} {Exciton relaxation cascade in two-dimensional
  transition metal dichalcogenides},\ }\href
  {https://doi.org/10.1038/s41598-018-25906-7} {\bibfield  {journal} {\bibinfo
  {journal} {Scientific Reports}\ }\textbf {\bibinfo {volume} {8}},\ \bibinfo
  {pages} {8238} (\bibinfo {year} {2018})}\BibitemShut {NoStop}
\bibitem [{\citenamefont {Zhang}\ \emph {et~al.}(2025)\citenamefont {Zhang},
  \citenamefont {Hu}, \citenamefont {Perfetto},\ and\ \citenamefont
  {Stefanucci}}]{zhang_long-lived_2025}
  \BibitemOpen
  \bibfield  {author} {\bibinfo {author} {\bibfnamefont {Z.}~\bibnamefont
  {Zhang}}, \bibinfo {author} {\bibfnamefont {W.}~\bibnamefont {Hu}}, \bibinfo
  {author} {\bibfnamefont {E.}~\bibnamefont {Perfetto}},\ and\ \bibinfo
  {author} {\bibfnamefont {G.}~\bibnamefont {Stefanucci}},\ }\bibfield  {title}
  {\bibinfo {title} {Long-lived coherence between incoherent excitons revealed
  by time-resolved angle-resolved photoemission spectroscopy: An exact
  solution},\ }\href {https://doi.org/10.1103/9spx-3hrq} {\bibfield  {journal}
  {\bibinfo  {journal} {Phys. Rev. B}\ }\textbf {\bibinfo {volume} {111}},\
  \bibinfo {pages} {235124} (\bibinfo {year} {2025})}\BibitemShut {NoStop}
\bibitem [{\citenamefont {Chen}\ \emph {et~al.}(2022)\citenamefont {Chen},
  \citenamefont {Sangalli},\ and\ \citenamefont
  {Bernardi}}]{chen_first-principles_2022}
  \BibitemOpen
  \bibfield  {author} {\bibinfo {author} {\bibfnamefont {H.-Y.}\ \bibnamefont
  {Chen}}, \bibinfo {author} {\bibfnamefont {D.}~\bibnamefont {Sangalli}},\
  and\ \bibinfo {author} {\bibfnamefont {M.}~\bibnamefont {Bernardi}},\
  }\bibfield  {title} {\bibinfo {title} {First-principles ultrafast exciton
  dynamics and time-domain spectroscopies: Dark-exciton mediated valley
  depolarization in monolayer {WS}e$_{2}$},\ }\href
  {https://doi.org/10.1103/PhysRevResearch.4.043203} {\bibfield  {journal}
  {\bibinfo  {journal} {Phys. Rev. Res.}\ }\textbf {\bibinfo {volume} {4}},\
  \bibinfo {pages} {043203} (\bibinfo {year} {2022})}\BibitemShut {NoStop}
\bibitem [{\citenamefont {Kremp}\ \emph {et~al.}(1984)\citenamefont {Kremp},
  \citenamefont {Kraeft},\ and\ \citenamefont {Lambert}}]{kremp_equation_1984}
  \BibitemOpen
  \bibfield  {author} {\bibinfo {author} {\bibfnamefont {D.}~\bibnamefont
  {Kremp}}, \bibinfo {author} {\bibfnamefont {W.}~\bibnamefont {Kraeft}},\ and\
  \bibinfo {author} {\bibfnamefont {A.}~\bibnamefont {Lambert}},\ }\bibfield
  {title} {\bibinfo {title} {Equation of state and ionization equilibrium for
  nonideal plasmas},\ }\href
  {https://doi.org/https://doi.org/10.1016/0378-4371(84)90120-1} {\bibfield
  {journal} {\bibinfo  {journal} {Physica A: Statistical Mechanics and its
  Applications}\ }\textbf {\bibinfo {volume} {127}},\ \bibinfo {pages} {72}
  (\bibinfo {year} {1984})}\BibitemShut {NoStop}
\bibitem [{\citenamefont {Semkat}\ \emph {et~al.}(2009)\citenamefont {Semkat},
  \citenamefont {Richter}, \citenamefont {Kremp}, \citenamefont {Manzke},
  \citenamefont {Kraeft},\ and\ \citenamefont
  {Henneberger}}]{semkat_ionization_2009}
  \BibitemOpen
  \bibfield  {author} {\bibinfo {author} {\bibfnamefont {D.}~\bibnamefont
  {Semkat}}, \bibinfo {author} {\bibfnamefont {F.}~\bibnamefont {Richter}},
  \bibinfo {author} {\bibfnamefont {D.}~\bibnamefont {Kremp}}, \bibinfo
  {author} {\bibfnamefont {G.}~\bibnamefont {Manzke}}, \bibinfo {author}
  {\bibfnamefont {W.-D.}\ \bibnamefont {Kraeft}},\ and\ \bibinfo {author}
  {\bibfnamefont {K.}~\bibnamefont {Henneberger}},\ }\bibfield  {title}
  {\bibinfo {title} {Ionization equilibrium in an excited semiconductor: Mott
  transition versus bose-einstein condensation},\ }\href
  {https://link.aps.org/doi/10.1103/PhysRevB.80.155201} {\bibfield  {journal}
  {\bibinfo  {journal} {Phys. Rev. B}\ }\textbf {\bibinfo {volume} {80}},\
  \bibinfo {pages} {155201} (\bibinfo {year} {2009})}\BibitemShut {NoStop}
\bibitem [{\citenamefont {Strinati}(1988)}]{strinati_application_1988}
  \BibitemOpen
  \bibfield  {author} {\bibinfo {author} {\bibfnamefont {G.}~\bibnamefont
  {Strinati}},\ }\bibfield  {title} {\bibinfo {title} {Application of the
  green's functions method to the study of the optical properties of
  semiconductors},\ }\href {https://doi.org/10.1007/BF02725962} {\bibfield
  {journal} {\bibinfo  {journal} {La Rivista del Nuovo Cimento (1978-1999)}\
  }\textbf {\bibinfo {volume} {11}},\ \bibinfo {pages} {1} (\bibinfo {year}
  {1988})}\BibitemShut {NoStop}
\bibitem [{not()}]{notepolaronic}
  \BibitemOpen
  \href@noop {} {}\bibinfo {note} {A tiny difference may arise from the
  polaronic-like shift of the exciton energy in the incoherent regime, see
  Ref.~\cite{stefanucci_from-carriers_2021}.}\BibitemShut {Stop}
\bibitem [{\citenamefont {Lin}\ \emph {et~al.}(2021)\citenamefont {Lin},
  \citenamefont {Ong}, \citenamefont {Bange}, \citenamefont {Faria~Junior},
  \citenamefont {Peng}, \citenamefont {Ziegler}, \citenamefont {Zipfel},
  \citenamefont {B{\"a}uml}, \citenamefont {Paradiso}, \citenamefont
  {Watanabe}, \citenamefont {Taniguchi}, \citenamefont {Strunk}, \citenamefont
  {Monserrat}, \citenamefont {Fabian}, \citenamefont {Chernikov}, \citenamefont
  {Qiu}, \citenamefont {Louie},\ and\ \citenamefont
  {Lupton}}]{lin_narrow-band_2021}
  \BibitemOpen
  \bibfield  {author} {\bibinfo {author} {\bibfnamefont {K.-Q.}\ \bibnamefont
  {Lin}}, \bibinfo {author} {\bibfnamefont {C.~S.}\ \bibnamefont {Ong}},
  \bibinfo {author} {\bibfnamefont {S.}~\bibnamefont {Bange}}, \bibinfo
  {author} {\bibfnamefont {P.~E.}\ \bibnamefont {Faria~Junior}}, \bibinfo
  {author} {\bibfnamefont {B.}~\bibnamefont {Peng}}, \bibinfo {author}
  {\bibfnamefont {J.~D.}\ \bibnamefont {Ziegler}}, \bibinfo {author}
  {\bibfnamefont {J.}~\bibnamefont {Zipfel}}, \bibinfo {author} {\bibfnamefont
  {C.}~\bibnamefont {B{\"a}uml}}, \bibinfo {author} {\bibfnamefont
  {N.}~\bibnamefont {Paradiso}}, \bibinfo {author} {\bibfnamefont
  {K.}~\bibnamefont {Watanabe}}, \bibinfo {author} {\bibfnamefont
  {T.}~\bibnamefont {Taniguchi}}, \bibinfo {author} {\bibfnamefont
  {C.}~\bibnamefont {Strunk}}, \bibinfo {author} {\bibfnamefont
  {B.}~\bibnamefont {Monserrat}}, \bibinfo {author} {\bibfnamefont
  {J.}~\bibnamefont {Fabian}}, \bibinfo {author} {\bibfnamefont
  {A.}~\bibnamefont {Chernikov}}, \bibinfo {author} {\bibfnamefont {D.~Y.}\
  \bibnamefont {Qiu}}, \bibinfo {author} {\bibfnamefont {S.~G.}\ \bibnamefont
  {Louie}},\ and\ \bibinfo {author} {\bibfnamefont {J.~M.}\ \bibnamefont
  {Lupton}},\ }\bibfield  {title} {\bibinfo {title} {Narrow-band high-lying
  excitons with negative-mass electrons in monolayer wse2},\ }\href
  {https://doi.org/10.1038/s41467-021-25499-2} {\bibfield  {journal} {\bibinfo
  {journal} {Nature Communications}\ }\textbf {\bibinfo {volume} {12}},\
  \bibinfo {pages} {5500} (\bibinfo {year} {2021})}\BibitemShut {NoStop}
\bibitem [{\citenamefont {Deilmann}\ and\ \citenamefont
  {Thygesen}(2019)}]{deilmann_finite-momentum_2019}
  \BibitemOpen
  \bibfield  {author} {\bibinfo {author} {\bibfnamefont {T.}~\bibnamefont
  {Deilmann}}\ and\ \bibinfo {author} {\bibfnamefont {K.~S.}\ \bibnamefont
  {Thygesen}},\ }\bibfield  {title} {\bibinfo {title} {Finite-momentum exciton
  landscape in mono- and bilayer transition metal dichalcogenides},\ }\href
  {https://doi.org/10.1088/2053-1583/ab0e1d} {\bibfield  {journal} {\bibinfo
  {journal} {2D Materials}\ }\textbf {\bibinfo {volume} {6}},\ \bibinfo {pages}
  {035003} (\bibinfo {year} {2019})}\BibitemShut {NoStop}
\bibitem [{\citenamefont {Foster}\ \emph {et~al.}(2014)\citenamefont {Foster},
  \citenamefont {Gurarie}, \citenamefont {Dzero},\ and\ \citenamefont
  {Yuzbashyan}}]{foster_quench-induced_2014}
  \BibitemOpen
  \bibfield  {author} {\bibinfo {author} {\bibfnamefont {M.~S.}\ \bibnamefont
  {Foster}}, \bibinfo {author} {\bibfnamefont {V.}~\bibnamefont {Gurarie}},
  \bibinfo {author} {\bibfnamefont {M.}~\bibnamefont {Dzero}},\ and\ \bibinfo
  {author} {\bibfnamefont {E.~A.}\ \bibnamefont {Yuzbashyan}},\ }\bibfield
  {title} {\bibinfo {title} {Quench-induced floquet topological $p$-wave
  superfluids},\ }\href
  {https://link.aps.org/doi/10.1103/PhysRevLett.113.076403} {\bibfield
  {journal} {\bibinfo  {journal} {Phys. Rev. Lett.}\ }\textbf {\bibinfo
  {volume} {113}},\ \bibinfo {pages} {076403} (\bibinfo {year}
  {2014})}\BibitemShut {NoStop}
\bibitem [{\citenamefont {Rustagi}\ and\ \citenamefont
  {Kemper}(2019)}]{rustagi_coherent_2019}
  \BibitemOpen
  \bibfield  {author} {\bibinfo {author} {\bibfnamefont {A.}~\bibnamefont
  {Rustagi}}\ and\ \bibinfo {author} {\bibfnamefont {A.~F.}\ \bibnamefont
  {Kemper}},\ }\bibfield  {title} {\bibinfo {title} {Coherent excitonic quantum
  beats in time-resolved photoemission measurements},\ }\href
  {https://doi.org/10.1103/PhysRevB.99.125303} {\bibfield  {journal} {\bibinfo
  {journal} {Phys. Rev. B}\ }\textbf {\bibinfo {volume} {99}},\ \bibinfo
  {pages} {125303} (\bibinfo {year} {2019})}\BibitemShut {NoStop}
\bibitem [{\citenamefont {Liu}\ \emph {et~al.}(2013)\citenamefont {Liu},
  \citenamefont {Shan}, \citenamefont {Yao}, \citenamefont {Yao},\ and\
  \citenamefont {Xiao}}]{PhysRevB.88.085433}
  \BibitemOpen
  \bibfield  {author} {\bibinfo {author} {\bibfnamefont {G.-B.}\ \bibnamefont
  {Liu}}, \bibinfo {author} {\bibfnamefont {W.-Y.}\ \bibnamefont {Shan}},
  \bibinfo {author} {\bibfnamefont {Y.}~\bibnamefont {Yao}}, \bibinfo {author}
  {\bibfnamefont {W.}~\bibnamefont {Yao}},\ and\ \bibinfo {author}
  {\bibfnamefont {D.}~\bibnamefont {Xiao}},\ }\bibfield  {title} {\bibinfo
  {title} {Three-band tight-binding model for monolayers of group-vib
  transition metal dichalcogenides},\ }\href
  {https://doi.org/10.1103/PhysRevB.88.085433} {\bibfield  {journal} {\bibinfo
  {journal} {Phys. Rev. B}\ }\textbf {\bibinfo {volume} {88}},\ \bibinfo
  {pages} {085433} (\bibinfo {year} {2013})}\BibitemShut {NoStop}
\bibitem [{\citenamefont {Chernikov}\ \emph {et~al.}(2014)\citenamefont
  {Chernikov}, \citenamefont {Berkelbach}, \citenamefont {Hill}, \citenamefont
  {Rigosi}, \citenamefont {Li}, \citenamefont {Aslan}, \citenamefont
  {Reichman}, \citenamefont {Hybertsen},\ and\ \citenamefont
  {Heinz}}]{PhysRevLett.113.076802}
  \BibitemOpen
  \bibfield  {author} {\bibinfo {author} {\bibfnamefont {A.}~\bibnamefont
  {Chernikov}}, \bibinfo {author} {\bibfnamefont {T.~C.}\ \bibnamefont
  {Berkelbach}}, \bibinfo {author} {\bibfnamefont {H.~M.}\ \bibnamefont
  {Hill}}, \bibinfo {author} {\bibfnamefont {A.}~\bibnamefont {Rigosi}},
  \bibinfo {author} {\bibfnamefont {Y.}~\bibnamefont {Li}}, \bibinfo {author}
  {\bibfnamefont {B.}~\bibnamefont {Aslan}}, \bibinfo {author} {\bibfnamefont
  {D.~R.}\ \bibnamefont {Reichman}}, \bibinfo {author} {\bibfnamefont {M.~S.}\
  \bibnamefont {Hybertsen}},\ and\ \bibinfo {author} {\bibfnamefont {T.~F.}\
  \bibnamefont {Heinz}},\ }\bibfield  {title} {\bibinfo {title} {Exciton
  binding energy and nonhydrogenic rydberg series in monolayer
  ${\mathrm{ws}}_{2}$},\ }\href
  {https://doi.org/10.1103/PhysRevLett.113.076802} {\bibfield  {journal}
  {\bibinfo  {journal} {Phys. Rev. Lett.}\ }\textbf {\bibinfo {volume} {113}},\
  \bibinfo {pages} {076802} (\bibinfo {year} {2014})}\BibitemShut {NoStop}
\bibitem [{\citenamefont {Hsu}\ \emph {et~al.}(2019)\citenamefont {Hsu},
  \citenamefont {Quan}, \citenamefont {Wang}, \citenamefont {Lu}, \citenamefont
  {Campbell}, \citenamefont {Chang}, \citenamefont {Li}, \citenamefont {Li},\
  and\ \citenamefont {Shih}}]{Hsu_2019}
  \BibitemOpen
  \bibfield  {author} {\bibinfo {author} {\bibfnamefont {W.-T.}\ \bibnamefont
  {Hsu}}, \bibinfo {author} {\bibfnamefont {J.}~\bibnamefont {Quan}}, \bibinfo
  {author} {\bibfnamefont {C.-Y.}\ \bibnamefont {Wang}}, \bibinfo {author}
  {\bibfnamefont {L.-S.}\ \bibnamefont {Lu}}, \bibinfo {author} {\bibfnamefont
  {M.}~\bibnamefont {Campbell}}, \bibinfo {author} {\bibfnamefont {W.-H.}\
  \bibnamefont {Chang}}, \bibinfo {author} {\bibfnamefont {L.-J.}\ \bibnamefont
  {Li}}, \bibinfo {author} {\bibfnamefont {X.}~\bibnamefont {Li}},\ and\
  \bibinfo {author} {\bibfnamefont {C.-K.}\ \bibnamefont {Shih}},\ }\bibfield
  {title} {\bibinfo {title} {Dielectric impact on exciton binding energy and
  quasiparticle bandgap in monolayer ws2 and wse2},\ }\href
  {https://doi.org/10.1088/2053-1583/ab072a} {\bibfield  {journal} {\bibinfo
  {journal} {2D Materials}\ }\textbf {\bibinfo {volume} {6}},\ \bibinfo {pages}
  {025028} (\bibinfo {year} {2019})}\BibitemShut {NoStop}
\bibitem [{\citenamefont {{Keldysh}}(1979)}]{keldysh}
  \BibitemOpen
  \bibfield  {author} {\bibinfo {author} {\bibfnamefont {L.~V.}\ \bibnamefont
  {{Keldysh}}},\ }\bibfield  {title} {\bibinfo {title} {{Coulomb interaction in
  thin semiconductor and semimetal films}},\ }\href@noop {} {\bibfield
  {journal} {\bibinfo  {journal} {Soviet Journal of Experimental and
  Theoretical Physics Letters}\ }\textbf {\bibinfo {volume} {29}},\ \bibinfo
  {pages} {658} (\bibinfo {year} {1979})}\BibitemShut {NoStop}
\bibitem [{\citenamefont {Cudazzo}\ \emph {et~al.}(2011)\citenamefont
  {Cudazzo}, \citenamefont {Tokatly},\ and\ \citenamefont
  {Rubio}}]{PhysRevB.84.085406}
  \BibitemOpen
  \bibfield  {author} {\bibinfo {author} {\bibfnamefont {P.}~\bibnamefont
  {Cudazzo}}, \bibinfo {author} {\bibfnamefont {I.~V.}\ \bibnamefont
  {Tokatly}},\ and\ \bibinfo {author} {\bibfnamefont {A.}~\bibnamefont
  {Rubio}},\ }\bibfield  {title} {\bibinfo {title} {Dielectric screening in
  two-dimensional insulators: Implications for excitonic and impurity states in
  graphane},\ }\href {https://doi.org/10.1103/PhysRevB.84.085406} {\bibfield
  {journal} {\bibinfo  {journal} {Phys. Rev. B}\ }\textbf {\bibinfo {volume}
  {84}},\ \bibinfo {pages} {085406} (\bibinfo {year} {2011})}\BibitemShut
  {NoStop}
\bibitem [{\citenamefont {Stier}\ \emph {et~al.}(2018)\citenamefont {Stier},
  \citenamefont {Wilson}, \citenamefont {Velizhanin}, \citenamefont {Kono},
  \citenamefont {Xu},\ and\ \citenamefont
  {Crooker}}]{stier_magnetooptics_2018}
  \BibitemOpen
  \bibfield  {author} {\bibinfo {author} {\bibfnamefont {A.~V.}\ \bibnamefont
  {Stier}}, \bibinfo {author} {\bibfnamefont {N.~P.}\ \bibnamefont {Wilson}},
  \bibinfo {author} {\bibfnamefont {K.~A.}\ \bibnamefont {Velizhanin}},
  \bibinfo {author} {\bibfnamefont {J.}~\bibnamefont {Kono}}, \bibinfo {author}
  {\bibfnamefont {X.}~\bibnamefont {Xu}},\ and\ \bibinfo {author}
  {\bibfnamefont {S.~A.}\ \bibnamefont {Crooker}},\ }\bibfield  {title}
  {\bibinfo {title} {Magnetooptics of exciton rydberg states in a monolayer
  semiconductor},\ }\href {https://doi.org/10.1103/PhysRevLett.120.057405}
  {\bibfield  {journal} {\bibinfo  {journal} {Phys. Rev. Lett.}\ }\textbf
  {\bibinfo {volume} {120}},\ \bibinfo {pages} {057405} (\bibinfo {year}
  {2018})}\BibitemShut {NoStop}
\bibitem [{\citenamefont {H\"user}\ \emph {et~al.}(2013)\citenamefont
  {H\"user}, \citenamefont {Olsen},\ and\ \citenamefont
  {Thygesen}}]{PhysRevB.88.245309}
  \BibitemOpen
  \bibfield  {author} {\bibinfo {author} {\bibfnamefont {F.}~\bibnamefont
  {H\"user}}, \bibinfo {author} {\bibfnamefont {T.}~\bibnamefont {Olsen}},\
  and\ \bibinfo {author} {\bibfnamefont {K.~S.}\ \bibnamefont {Thygesen}},\
  }\bibfield  {title} {\bibinfo {title} {How dielectric screening in
  two-dimensional crystals affects the convergence of excited-state
  calculations: Monolayer mos${}_{2}$},\ }\href
  {https://doi.org/10.1103/PhysRevB.88.245309} {\bibfield  {journal} {\bibinfo
  {journal} {Phys. Rev. B}\ }\textbf {\bibinfo {volume} {88}},\ \bibinfo
  {pages} {245309} (\bibinfo {year} {2013})}\BibitemShut {NoStop}
\bibitem [{\citenamefont {Ridolfi}\ \emph {et~al.}(2018)\citenamefont
  {Ridolfi}, \citenamefont {Lewenkopf},\ and\ \citenamefont
  {Pereira}}]{PhysRevB.97.205409}
  \BibitemOpen
  \bibfield  {author} {\bibinfo {author} {\bibfnamefont {E.}~\bibnamefont
  {Ridolfi}}, \bibinfo {author} {\bibfnamefont {C.~H.}\ \bibnamefont
  {Lewenkopf}},\ and\ \bibinfo {author} {\bibfnamefont {V.~M.}\ \bibnamefont
  {Pereira}},\ }\bibfield  {title} {\bibinfo {title} {Excitonic structure of
  the optical conductivity in ${\mathrm{mos}}_{2}$ monolayers},\ }\href
  {https://doi.org/10.1103/PhysRevB.97.205409} {\bibfield  {journal} {\bibinfo
  {journal} {Phys. Rev. B}\ }\textbf {\bibinfo {volume} {97}},\ \bibinfo
  {pages} {205409} (\bibinfo {year} {2018})}\BibitemShut {NoStop}
\bibitem [{\citenamefont {Groenewald}\ \emph {et~al.}(2016)\citenamefont
  {Groenewald}, \citenamefont {R\"osner}, \citenamefont {Sch\"onhoff},
  \citenamefont {Haas},\ and\ \citenamefont {Wehling}}]{PhysRevB.93.205145}
  \BibitemOpen
  \bibfield  {author} {\bibinfo {author} {\bibfnamefont {R.~E.}\ \bibnamefont
  {Groenewald}}, \bibinfo {author} {\bibfnamefont {M.}~\bibnamefont
  {R\"osner}}, \bibinfo {author} {\bibfnamefont {G.}~\bibnamefont
  {Sch\"onhoff}}, \bibinfo {author} {\bibfnamefont {S.}~\bibnamefont {Haas}},\
  and\ \bibinfo {author} {\bibfnamefont {T.~O.}\ \bibnamefont {Wehling}},\
  }\bibfield  {title} {\bibinfo {title} {Valley plasmonics in transition metal
  dichalcogenides},\ }\href {https://doi.org/10.1103/PhysRevB.93.205145}
  {\bibfield  {journal} {\bibinfo  {journal} {Phys. Rev. B}\ }\textbf {\bibinfo
  {volume} {93}},\ \bibinfo {pages} {205145} (\bibinfo {year}
  {2016})}\BibitemShut {NoStop}
\bibitem [{\citenamefont {Steinhoff}\ \emph {et~al.}(2014)\citenamefont
  {Steinhoff}, \citenamefont {R{\"o}sner}, \citenamefont {Jahnke},
  \citenamefont {Wehling},\ and\ \citenamefont {Gies}}]{steinhoff2}
  \BibitemOpen
  \bibfield  {author} {\bibinfo {author} {\bibfnamefont {A.}~\bibnamefont
  {Steinhoff}}, \bibinfo {author} {\bibfnamefont {M.}~\bibnamefont
  {R{\"o}sner}}, \bibinfo {author} {\bibfnamefont {F.}~\bibnamefont {Jahnke}},
  \bibinfo {author} {\bibfnamefont {T.~O.}\ \bibnamefont {Wehling}},\ and\
  \bibinfo {author} {\bibfnamefont {C.}~\bibnamefont {Gies}},\ }\bibfield
  {title} {\bibinfo {title} {Influence of excited carriers on the optical and
  electronic properties of mos2},\ }\href {https://doi.org/10.1021/nl500595u}
  {\bibfield  {journal} {\bibinfo  {journal} {Nano Letters}\ }\textbf {\bibinfo
  {volume} {14}},\ \bibinfo {pages} {3743} (\bibinfo {year}
  {2014})}\BibitemShut {NoStop}
\bibitem [{\citenamefont {Berkelbach}\ \emph {et~al.}(2013)\citenamefont
  {Berkelbach}, \citenamefont {Hybertsen},\ and\ \citenamefont
  {Reichman}}]{PhysRevB.88.045318}
  \BibitemOpen
  \bibfield  {author} {\bibinfo {author} {\bibfnamefont {T.~C.}\ \bibnamefont
  {Berkelbach}}, \bibinfo {author} {\bibfnamefont {M.~S.}\ \bibnamefont
  {Hybertsen}},\ and\ \bibinfo {author} {\bibfnamefont {D.~R.}\ \bibnamefont
  {Reichman}},\ }\bibfield  {title} {\bibinfo {title} {Theory of neutral and
  charged excitons in monolayer transition metal dichalcogenides},\ }\href
  {https://doi.org/10.1103/PhysRevB.88.045318} {\bibfield  {journal} {\bibinfo
  {journal} {Phys. Rev. B}\ }\textbf {\bibinfo {volume} {88}},\ \bibinfo
  {pages} {045318} (\bibinfo {year} {2013})}\BibitemShut {NoStop}
\end{thebibliography}


\end{document}
===== END FILE: TR-ARPES-X.tex =====

===== BEGIN FILE: ourmacros.tex =====

\newcommand{\be}{\begin{equation}}
\newcommand{\ee}{\end{equation}}
\newcommand{\bea}{\begin{eqnarray}}
\newcommand{\eea}{\end{eqnarray}}
\newcommand{\bi}{\bibitem}
\newcommand\diag[2]{
  \ncline[linewidth=1pt,nodesep=-30pt]{#1}{#2}}


\def\a{\alpha}
\def\b{\beta}
\def\g{\gamma}
\def\G{\Gamma}
\def\d{\delta}
\def\D{\Delta}
\def\e{\epsilon}
\def\ve{\varepsilon}
\def\h{\eta}
\def\th{\theta}
\def\Th{\Theta}
\def\vt{\vartheta}
\def\k{\kappa}
\def\l{\lambda}
\def\L{\Lambda}
\def\m{\mu}
\def\n{\nu}
\def\c{\xi}
\def\C{\Xi}
\def\p{\pi}
\def\P{\Pi}
\def\r{\rho}
\def\s{\sigma}
\def\S{\Sigma}
\def\t{\tau}
\def\f{\phi}
\def\vf{\varphi}
\def\F{\Phi}
\def\x{\chi}
\def\w{\omega}
\def\W{\Omega}
\def\q{\psi}
\def\Q{\Psi}
\def\z{\zeta}


\def\bga{\mbox{\boldmath $\alpha$}}
\def\bgb{\mbox{\boldmath $\beta$}}
\def\bgg{\mbox{\boldmath $\gamma$}}
\def\bgG{\mbox{\boldmath $\Gamma$}}
\def\bgd{\mbox{\boldmath $\delta$}}
\def\bgD{\mbox{\boldmath $\Delta$}}
\def\bge{\mbox{\boldmath $\epsilon$}}
\def\bgve{\mbox{\boldmath $\varepsilon$}}
\def\bgh{\mbox{\boldmath $\eta$}}
\def\bgth{\mbox{\boldmath $\theta$}}
\def\bgk{\mbox{\boldmath $\kappa$}}
\def\bgl{\mbox{\boldmath $\lambda$}}
\def\bgL{\mbox{\boldmath $\Lambda$}}
\def\bgm{\mbox{\boldmath $\mu$}}
\def\bgn{\mbox{\boldmath $\nu$}}
\def\bgc{\mbox{\boldmath $\xi$}}
\def\bgC{\mbox{\boldmath $\Xi$}}
\def\bgp{\mbox{\boldmath $\pi$}}
\def\bgP{\mbox{\boldmath $\Pi$}}
\def\bgr{\mbox{\boldmath $\rho$}}
\def\bgs{\mbox{\boldmath $\sigma$}}
\def\bgS{\mbox{\boldmath $\Sigma$}}
\def\bgt{\mbox{\boldmath $\tau$}}
\def\bgf{\mbox{\boldmath $\phi$}}
\def\bgvf{\mbox{\boldmath $\varphi$}}
\def\bgF{\mbox{\boldmath $\Phi$}}
\def\bgx{\mbox{\boldmath $\chi$}}
\def\bgw{\mbox{\boldmath $\omega$}}
\def\bgW{\mbox{\boldmath $\Omega$}}
\def\bgq{\mbox{\boldmath $\psi$}}
\def\bgQ{\mbox{\boldmath $\Psi$}}
\def\bgz{\mbox{\boldmath $\zeta$}}


\def\calgG{\mbox{$\mathit{\Gamma}$}}
\def\calgD{\mbox{$\mathit{\Delta}$}}
\def\calgL{\mbox{$\mathit{\Lambda}$}}
\def\calgC{\mbox{$\mathit{\Xi}$}}
\def\calgP{\mbox{$\mathit{\Pi}$}}
\def\calgS{\mbox{$\mathit{\Sigma}$}}
\def\calgF{\mbox{$\mathit{\Phi}$}}
\def\calgW{\mbox{$\mathit{\Omega}$}}
\def\calgQ{\mbox{$\mathit{\Psi}$}}


\def\bcalgG{\mbox{\boldmath $\mathit{\Gamma}$}}
\def\bcalgD{\mbox{\boldmath $\mathit{\Delta}$}}
\def\bcalgL{\mbox{\boldmath $\mathit{\Lambda}$}}
\def\bcalgC{\mbox{\boldmath $\mathit{\Xi}$}}
\def\bcalgP{\mbox{\boldmath $\mathit{\Pi}$}}
\def\bcalgS{\mbox{\boldmath $\mathit{\Sigma}$}}
\def\bcalgF{\mbox{\boldmath $\mathit{\Phi}$}}
\def\bcalgW{\mbox{\boldmath $\mathit{\Omega}$}}
\def\bcalgQ{\mbox{\boldmath $\mathit{\Psi}$}}


\def\bla{{\mathbf a}}
\def\blb{{\mathbf b}}
\def\blc{{\mathbf c}}
\def\bld{{\mathbf d}}
\def\ble{{\mathbf e}}
\def\blf{{\mathbf f}}
\def\blg{{\mathbf g}}
\def\blh{{\mathbf h}}
\def\bli{{\mathbf i}}
\def\blj{{\mathbf j}}
\def\blk{{\mathbf k}}
\def\bll{{\mathbf l}}
\def\blm{{\mathbf m}}
\def\bln{{\mathbf n}}
\def\blo{{\mathbf o}}
\def\blp{{\mathbf p}}
\def\blq{{\mathbf q}}
\def\blr{{\mathbf r}}
\def\bls{{\mathbf s}}
\def\blt{{\mathbf t}}
\def\blu{{\mathbf u}}
\def\blv{{\mathbf v}}
\def\blw{{\mathbf w}}
\def\blx{{\mathbf x}}
\def\bly{{\mathbf y}}
\def\blz{{\mathbf z}}


\def\bcalla{\mbox{\boldmath $a$}}
\def\bcallb{\mbox{\boldmath $b$}}
\def\bcallc{\mbox{\boldmath $c$}}
\def\bcalld{\mbox{\boldmath $d$}}
\def\bcalle{\mbox{\boldmath $e$}}
\def\bcallf{\mbox{\boldmath $f$}}
\def\bcallg{\mbox{\boldmath $g$}}
\def\bcallh{\mbox{\boldmath $h$}}
\def\bcalli{\mbox{\boldmath $i$}}
\def\bcallj{\mbox{\boldmath $j$}}
\def\bcallk{\mbox{\boldmath $k$}}
\def\bcalll{\mbox{\boldmath $l$}}
\def\bcallm{\mbox{\boldmath $m$}}
\def\bcalln{\mbox{\boldmath $n$}}
\def\bcallo{\mbox{\boldmath $o$}}
\def\bcallp{\mbox{\boldmath $p$}}
\def\bcallq{\mbox{\boldmath $q$}}
\def\bcallr{\mbox{\boldmath $r$}}
\def\bcalls{\mbox{\boldmath $s$}}
\def\bcallt{\mbox{\boldmath $t$}}
\def\bcallu{\mbox{\boldmath $u$}}
\def\bcallv{\mbox{\boldmath $v$}}
\def\bcallw{\mbox{\boldmath $w$}}
\def\bcallx{\mbox{\boldmath $x$}}
\def\bcally{\mbox{\boldmath $y$}}
\def\bcallz{\mbox{\boldmath $z$}}


\def\blA{{\mathbf A}}
\def\blB{{\mathbf B}}
\def\blC{{\mathbf C}}
\def\blD{{\mathbf D}}
\def\blE{{\mathbf E}}
\def\blF{{\mathbf F}}
\def\blG{{\mathbf G}}
\def\blH{{\mathbf H}}
\def\blI{{\mathbf I}}
\def\blJ{{\mathbf J}}
\def\blK{{\mathbf K}}
\def\blL{{\mathbf L}}
\def\blM{{\mathbf M}}
\def\blN{{\mathbf N}}
\def\blO{{\mathbf O}}
\def\blP{{\mathbf P}}
\def\blQ{{\mathbf Q}}
\def\blR{{\mathbf R}}
\def\blS{{\mathbf S}}
\def\blT{{\mathbf T}}
\def\blU{{\mathbf U}}
\def\blV{{\mathbf V}}
\def\blW{{\mathbf W}}
\def\blX{{\mathbf X}}
\def\blY{{\mathbf Y}}
\def\blZ{{\mathbf Z}}


\def\bilA{\mbox{\boldmath $A$}}
\def\bilB{\mbox{\boldmath $B$}}
\def\bilC{\mbox{\boldmath $C$}}
\def\bilD{\mbox{\boldmath $D$}}
\def\bilE{\mbox{\boldmath $E$}}
\def\bilF{\mbox{\boldmath $F$}}
\def\bilG{\mbox{\boldmath $G$}}
\def\bilH{\mbox{\boldmath $H$}}
\def\bilI{\mbox{\boldmath $I$}}
\def\bilJ{\mbox{\boldmath $J$}}
\def\bilK{\mbox{\boldmath $K$}}
\def\bilL{\mbox{\boldmath $L$}}
\def\bilM{\mbox{\boldmath $M$}}
\def\bilN{\mbox{\boldmath $N$}}
\def\bilO{\mbox{\boldmath $O$}}
\def\bilP{\mbox{\boldmath $P$}}
\def\bilQ{\mbox{\boldmath $Q$}}
\def\bilR{\mbox{\boldmath $R$}}
\def\bilS{\mbox{\boldmath $S$}}
\def\bilT{\mbox{\boldmath $T$}}
\def\bilU{\mbox{\boldmath $U$}}
\def\bilV{\mbox{\boldmath $V$}}
\def\bilW{\mbox{\boldmath $W$}}
\def\bilX{\mbox{\boldmath $X$}}
\def\bilY{\mbox{\boldmath $Y$}}
\def\bilZ{\mbox{\boldmath $Z$}}


\def\callA{\mbox{$\mathcal{A}$}}
\def\callB{\mbox{$\mathcal{B}$}}
\def\callC{\mbox{$\mathcal{C}$}}
\def\callD{\mbox{$\mathcal{D}$}}
\def\callE{\mbox{$\mathcal{E}$}}
\def\callF{\mbox{$\mathcal{F}$}}
\def\callG{\mbox{$\mathcal{G}$}}
\def\callH{\mbox{$\mathcal{H}$}}
\def\callI{\mbox{$\mathcal{I}$}}
\def\callJ{\mbox{$\mathcal{J}$}}
\def\callK{\mbox{$\mathcal{K}$}}
\def\callL{\mbox{$\mathcal{L}$}}
\def\callM{\mbox{$\mathcal{M}$}}
\def\callN{\mbox{$\mathcal{N}$}}
\def\callO{\mbox{$\mathcal{O}$}}
\def\callP{\mbox{$\mathcal{P}$}}
\def\callQ{\mbox{$\mathcal{Q}$}}
\def\callR{\mbox{$\mathcal{R}$}}
\def\callS{\mbox{$\mathcal{S}$}}
\def\callT{\mbox{$\mathcal{T}$}}
\def\callU{\mbox{$\mathcal{U}$}}
\def\callV{\mbox{$\mathcal{V}$}}
\def\callW{\mbox{$\mathcal{W}$}}
\def\callX{\mbox{$\mathcal{X}$}}
\def\callY{\mbox{$\mathcal{Y}$}}
\def\callZ{\mbox{$\mathcal{Z}$}}


\def\bcallA{\mbox{\boldmath $\mathcal{A}$}}
\def\bcallB{\mbox{\boldmath $\mathcal{B}$}}
\def\bcallC{\mbox{\boldmath $\mathcal{C}$}}
\def\bcallD{\mbox{\boldmath $\mathcal{D}$}}
\def\bcallE{\mbox{\boldmath $\mathcal{E}$}}
\def\bcallF{\mbox{\boldmath $\mathcal{F}$}}
\def\bcallG{\mbox{\boldmath $\mathcal{G}$}}
\def\bcallH{\mbox{\boldmath $\mathcal{H}$}}
\def\bcallI{\mbox{\boldmath $\mathcal{I}$}}
\def\bcallJ{\mbox{\boldmath $\mathcal{J}$}}
\def\bcallK{\mbox{\boldmath $\mathcal{K}$}}
\def\bcallL{\mbox{\boldmath $\mathcal{L}$}}
\def\bcallM{\mbox{\boldmath $\mathcal{M}$}}
\def\bcallN{\mbox{\boldmath $\mathcal{N}$}}
\def\bcallO{\mbox{\boldmath $\mathcal{O}$}}
\def\bcallP{\mbox{\boldmath $\mathcal{P}$}}
\def\bcallQ{\mbox{\boldmath $\mathcal{Q}$}}
\def\bcallR{\mbox{\boldmath $\mathcal{R}$}}
\def\bcallS{\mbox{\boldmath $\mathcal{S}$}}
\def\bcallT{\mbox{\boldmath $\mathcal{T}$}}
\def\bcallU{\mbox{\boldmath $\mathcal{U}$}}
\def\bcallV{\mbox{\boldmath $\mathcal{V}$}}
\def\bcallW{\mbox{\boldmath $\mathcal{W}$}}
\def\bcallX{\mbox{\boldmath $\mathcal{X}$}}
\def\bcallY{\mbox{\boldmath $\mathcal{Y}$}}
\def\bcallZ{\mbox{\boldmath $\mathcal{Z}$}}



\def\ua{\uparrow}
\def\da{\downarrow}
\def\ra{\rightarrow}
\def\la{\leftarrow}
\def\La{\Leftarrow}
\def\Ra{\Rightarrow}
\def\de{\partial}
\def\iif{\infty}
\def\bra{\langle}
\def\ket{\rangle}
\def\grad{\mbox{\boldmath $\nabla$}}
\def\minigrad{\mbox{{\footnotesize \boldmath $\nabla$}}}
\def\Tr{{\rm Tr}}
\def\Re{{\rm Re}}
\def\Im{{\rm Im}}




\def\iu{{\rm i}}
\def\1op{\hat{\mathbbm{1}}}
\def\nn{\nonumber}
\def\fbgc{\mbox{\scalebox{.8}{$\scriptscriptstyle \bgc$}}}
\def\bz{\mathbf 0}
===== END FILE: ourmacros.tex =====
